\documentclass{article}

\usepackage[a4paper, margin=1in]{geometry}
\usepackage{amssymb}
\usepackage{amsmath}
\usepackage{amsthm}
\usepackage{mathtools} % For \coloneqq
\usepackage{bm} % For \boldsymbol
\usepackage[dvipsnames]{xcolor} % Use dvipsnames for Maroon color if needed
\usepackage{enumitem}
\usepackage[authoryear]{natbib}
\usepackage{xr-hyper}
\usepackage{hyperref}
\usepackage{face-macros}

% Import labels from main.tex while ignoring BibTeX citation definitions in main.aux
% (otherwise natbib will warn about multiply-defined citations).
\makeatletter
\let\XR@old@bibcite\bibcite
\renewcommand\bibcite[2]{}
\externaldocument{main}
\let\bibcite\XR@old@bibcite
\makeatother

% --- Theorem Environments (from user's preamble) ---
\newtheorem{theorem}{Theorem}[section]
\newtheorem{lemma}[theorem]{Lemma}
\newtheorem{corollary}[theorem]{Corollary}
\newtheorem{assumption}{Assumption}[section]
\newtheorem{proposition}[theorem]{Proposition}
\newtheorem{definition}{Definition}[section]
\newtheorem{remark}{Remark}[section]

\newcommand{\hj}[1]{\textcolor{Maroon}{[Marquis: #1]}}

% --- Document Start ---
\begin{document}

\title{Appendix: Proofs of Asymptotic Properties for FACE-HD}
\author{Sinian Zhang}
\date{}
\maketitle

\appendix

\section{Asymptotic Normality of the Cross-Fitted Estimator}
\subsection{Setup and Notation}

We analyze the asymptotic properties of the $K_f$-fold cross-fitted estimator for a single target-source pair $(t, s)$, derived using the two-level cross-fitting strategy described in the main manuscript's appendix (Section~\ref{app:two_level_cf}).
Let $K_f$ be a fixed integer (e.g., 5 or 10). The data $\mathcal{D} = \mathcal{D}_t \cup \mathcal{D}_s$ is partitioned into $K_f$ folds, $\mathcal{D} = \cup_{k=1}^{K_f} \mathcal{D}_k$, where $\mathcal{D}_k = \mathcal{D}_{t,k} \cup \mathcal{D}_{s,k}$. Let $N_k = |\mathcal{D}_k| \approx N/K_f$.
For notational simplicity, we assume the folds are balanced within each site (e.g., $N_t$ and $N_s$ are divisible by $K_f$) so that $|\mathcal{D}_{t,k}|=N_t/K_f$ and $|\mathcal{D}_{s,k}|=N_s/K_f$ for all $k$; the minor rounding effects in practice are $o_p(N^{-1/2})$ and do not affect the first-order results.
For any function $f(D)$ of a single observation $D=(R,\boldsymbol{X},A,Y)$, we write $\tildeE_k[f(D)] \coloneqq \frac{1}{N_k}\sum_{i \in \mathcal{D}_k} f(D_i)$ for the empirical average over the combined $k$-th fold.
Let $\mathcal{D}^{(-k)} = \mathcal{D} \setminus \mathcal{D}_k$ denote the out-of-fold data used for training nuisance models evaluated on fold $k$. The crucial aspect of cross-fitting is that the data used for training, $\mathcal{D}^{(-k)}$, is independent of the data used for evaluation, $\mathcal{D}_k$.

Let $\vartheta = (\boldsymbol{\alpha}^1, \boldsymbol{\gamma}_1)$ be the vector of nuisance parameters. For this fixed pair $(t,s)$ we write $\boldsymbol{\alpha}^1 \equiv \boldsymbol{\alpha}_{t,s}^1$ and $\boldsymbol{\gamma}_1 \equiv \boldsymbol{\gamma}_{s,1}$. Let $\vartheta^* = (\boldsymbol{\alpha}^{1,*}, \boldsymbol{\gamma}_1^*)$ denote the population targets of the calibrated estimation step (Equations~\eqref{eq:gamma_calibrated_loss} and \eqref{eq:alpha_calibrated_loss} in the main manuscript), and let $\widehat{\vartheta}^{(-k)} = (\widehat{\boldsymbol{\alpha}}_{t,s}^{1,(-k)}, \widehat{\boldsymbol{\gamma}}_{s,1}^{(-k)})$ be the nuisance estimators trained on data $\mathcal{D}^{(-k)}$ using the two-level cross-fitting procedure.
Here the superscript $(-k)$ denotes the \emph{final} two-level cross-fitted nuisance estimators used for evaluating fold $k$ (including any inner-fold averaging), and hence $\widehat{\vartheta}^{(-k)}$ is independent of $\mathcal{D}_k$.

To ensure numerical stability and theoretical regularity (bounded weights), we define a truncation function $\mathcal{T}(x) = \text{sign}(x)\min(|x|, M_{\tau})$, where $M_{\tau}=M_{\tau,N}$ is allowed to depend on $N$ (notation suppressed). (We avoid the symbol $\tau(\cdot)$ to prevent collision with the ATE $\tau$ in the main manuscript.)
The truncated weighting function is defined as:
\[
w_\tau(\boldsymbol{X}; \boldsymbol{\gamma}) \coloneqq \frac{I(A=1)}{e^{\mathcal{T}(\phi(\boldsymbol{X})^T \boldsymbol{\gamma})}}.
\]
Define also the (untruncated) weight
\[
w(\boldsymbol{X}; \boldsymbol{\gamma}) \coloneqq \frac{I(A=1)}{e^{\phi(\boldsymbol{X})^T \boldsymbol{\gamma}}}.
\]
\textbf{Remark on Truncation / Clipping:} The truncation function $\mathcal{T}$ is applied \textit{only in the calibrated loss functions} (Equations~\eqref{eq:gamma_calibrated_loss} and \eqref{eq:alpha_calibrated_loss} in the main manuscript) to ensure numerical stability during nuisance estimation; the estimators and influence functions below use the untruncated weight $w(\boldsymbol{X};\boldsymbol{\gamma})$, consistent with the main manuscript. Under Assumption~\ref{asmp:overlap}(c), truncation is asymptotically inactive at the population targets (so $w_\tau(\boldsymbol{X};\boldsymbol{\gamma}_1^*) = w(\boldsymbol{X};\boldsymbol{\gamma}_1^*)$ and $\psi'(\mathcal{T}(\phi(\boldsymbol{X})^T\boldsymbol{\alpha}^{1,*}))=\psi'(\phi(\boldsymbol{X})^T\boldsymbol{\alpha}^{1,*})$ w.p.a.1), and hence the population first-order conditions induced by the calibrated losses are unchanged up to an $o(1)$ term.

In the software implementation, this corresponds to using a finite $M_{\tau}$ during nuisance optimization, while the default inference computations use $M_{\tau,\mathrm{inference}}=\infty$ (i.e., no truncation inside the final DR estimating equation). In addition, low-level C++ kernels retain very broad machine-safety bounds (e.g., \texttt{WEIGHT\_MIN/MAX}, \texttt{RATIO\_CLIP}) to prevent overflow/underflow; these are numerical safeguards and are treated as asymptotically inactive under the stated moment/overlap conditions. Setting a finite $M_{\tau,\mathrm{inference}}$ corresponds to the clipped estimating equation discussed below.

If, instead, one \emph{also} clips the weights in the final AIPW/DR estimating equation (i.e., replaces $w$ by $w_\tau$ inside $\widehat{\mu}^1_{t,s}$), then the resulting estimator has a different influence function in finite samples. Under Assumption~\ref{asmp:overlap}(c) (clipping probability $o(N^{-1})$ at the population targets), the difference between the clipped and unclipped estimators is $o_p(N^{-1/2})$ and the first-order limit theory is unchanged; otherwise, the clipping-induced bias must be accounted for (see, e.g., \citet{tan2020model}).

\subsection{Implementation Conventions Used in the Proof Mapping}
For clarity, we record the conventions used to align implementation details with the asymptotic statements:
\begin{itemize}
    \item \textbf{Fold-wise centering of IF components.} Variance and covariance components are computed from centered fold-level influence-function vectors, consistent with the sample-moment definitions in the main manuscript.
    \item \textbf{Weight-estimation variance vs final inference variance.} The penalized quadratic variance objective is used for estimating aggregation weights, whereas final cross-fitted uncertainty is computed from the sample-level unified estimating equation.
    \item \textbf{Finite-sample convexity stabilization.} In finite samples, the plug-in covariance matrix in the aggregation objective can be nearly indefinite; implementation-level symmetrization and minimal ridge stabilization are treated as numerical devices that vanish asymptotically under the regularity conditions.
\end{itemize}

\paragraph{Assumption-to-Implementation Map.}
Table~\ref{tab:assumption_code_map} provides a concise correspondence between core assumptions/claims and implementation controls.
\begin{table}[h]
\centering
\small
\begin{tabular}{|l|l|l|}
\hline
\textbf{Theoretical condition} & \textbf{Role in proof} & \textbf{Implementation anchor} \\
\hline
Fold-wise centering of IF components & Consistent variance/covariance moments & Fold-level centering in \texttt{R/cross\_fitting\_aggregation.R} \\
\hline
Training-only truncation ($M_{\tau}$) & Nuisance stability, same first-order target & Calibrated losses in \texttt{R/cross\_fitting\_algorithms.R} \\
\hline
Inference untruncated by default & Matches unclipped IF expansion & \texttt{M\_TAU\_INFERENCE\_DEFAULT = Inf} in \texttt{R/constants.R} \\
\hline
Finite-sample QP convexity handling & Well-posed optimization in practice & Symmetrize + PSD ridge in \texttt{src/weight\_optimization.cpp} \\
\hline
\end{tabular}
\caption{Assumption-to-implementation correspondence for reproducible inference.}
\label{tab:assumption_code_map}
\end{table}

We assume $\phi(\boldsymbol{X})$ includes an intercept term.
Let $p \coloneqq \dim(\phi(\boldsymbol{X}))$.
For canonical exponential-family GLMs, we use the standard notation that the negative log-likelihood takes the form $-Y\theta+\Psi(\theta)$ with linear predictor $\theta=\phi(\boldsymbol{X})^T\boldsymbol{\alpha}$, and that $\psi(\theta)=\Psi'(\theta)$; $\psi'(\theta)$ denotes the derivative with respect to $\theta$. We also use the shorthand $\psi(\phi(\boldsymbol{X});\boldsymbol{\alpha}) \equiv \psi(\phi(\boldsymbol{X})^T\boldsymbol{\alpha})$.
To avoid confusion with the notation used in the reference PDFs (in the \texttt{proof ref} folder), we consistently follow the main manuscript's notation: $\phi(\boldsymbol{X})$ denotes the feature/basis expansion and $\psi(\cdot)$ denotes the GLM conditional mean. (For example, some references use $\phi(\cdot)$ and $\Phi(\cdot)$ for a Normal density/CDF; we do \emph{not} use that convention here.)

As in the main manuscript, the target estimand is the target-site treated potential outcome mean $\mu^1 \equiv \E_t[Y(1)]$.
The $K_f$-fold cross-fitted estimator for $\mu^1$ is:
\begin{align*}
\widehat{\mu}^{1,\mathrm{CF}}_{t,s}
= \frac{1}{K_f}\sum_{k=1}^{K_f}
\Big\{
    \tildeE_{t,k}\big[\psi(\phi(\boldsymbol{X}); \widehat{\boldsymbol{\alpha}}_{t,s}^{1,(-k)})\big]
    + \tildeE_{s,k}\!\left[
        w(\boldsymbol{X};\widehat{\boldsymbol{\gamma}}_{s,1}^{(-k)})
        \big(Y - \psi(\phi(\boldsymbol{X}); \widehat{\boldsymbol{\alpha}}_{t,s}^{1,(-k)})\big)
    \right]
\Big\}.
\end{align*}
Here,
\[
\tildeE_{t,k}[g(D)] \coloneqq \frac{1}{N_{t,k}}\sum_{i \in \mathcal{D}_{t,k}} g(D_i),
\qquad
\tildeE_{s,k}[g(D)] \coloneqq \frac{1}{N_{s,k}}\sum_{i \in \mathcal{D}_{s,k}} g(D_i),
\]
denote sample averages over the $k$-th fold of the target and source data, with $N_{t,k} = |\mathcal{D}_{t,k}|$, $N_{s,k} = |\mathcal{D}_{s,k}|$, and $N_k = N_{t,k}+N_{s,k}$.
Let $N = N_t + N_s$ be the total sample size, and let $\rho_t = N_t/N$ and $\rho_s = N_s/N$.
We write $\E[\cdot]$ for expectation under the pooled distribution, i.e., $\E[g(D)] = \rho_t \E_t[g(D)] + \rho_s \E_s[g(D)]$.
Throughout the asymptotics we consider the standard two-sample (multi-sample) regime with sampling fractions bounded away from $0$ and $1$, i.e., $N_t/N\to\rho_t\in(0,1)$ and $N_s/N\to\rho_s\in(0,1)$. Under overlap/positivity for the relevant treatment arm, the effective sample sizes within each arm also scale linearly with $N$ (e.g., $\sum_{i\in\mathcal{D}_s}I(A_i=1)=\Theta_p(N_s)$), so writing high-dimensional rates in terms of $N$ is equivalent (up to constants) to writing them in terms of the arm-specific sample sizes.
Under the balanced-fold assumption above, the estimator can equivalently be written using the combined-fold empirical average as
\[
\widehat{\mu}^{1,\mathrm{CF}}_{t,s}
= \frac{1}{K_f}\sum_{k=1}^{K_f}\tildeE_k\!\left[\widetilde{\varphi}_{t,s}(D; \widehat{\vartheta}^{(-k)})\right],
\]
which matches the $\tildeE_k$ notation used in the remainder decomposition below.

Define the (uncentered) score contribution at a generic nuisance parameter $\vartheta = (\boldsymbol{\alpha}^1, \boldsymbol{\gamma}_1)$:
\begin{align*}
\widetilde{\varphi}_{t,s}(D_i; \vartheta)
&:= \frac{I(R_i=t)}{\rho_t}\psi(\phi(\boldsymbol{X}_i); \boldsymbol{\alpha}^1)
+ \frac{I(R_i=s)}{\rho_s} w(\boldsymbol{X}_i;\boldsymbol{\gamma}_1)\{Y_i - \psi(\phi(\boldsymbol{X}_i); \boldsymbol{\alpha}^1)\}.
\end{align*}
Define the corresponding population estimand at $\vartheta$ as
\begin{align*}
\mu^1(\vartheta)
&:= \E[\widetilde{\varphi}_{t,s}(D; \vartheta)]
= \E_t[\psi(\phi(\boldsymbol{X}); \boldsymbol{\alpha}^1)]
+ \E_s[w(\boldsymbol{X};\boldsymbol{\gamma}_1)\{Y - \psi(\phi(\boldsymbol{X}); \boldsymbol{\alpha}^1)\}].
\end{align*}


\subsection{Common Notation and Definitions}\label{sec:notation_defs}

To avoid defining objects inside later statements, we collect here the auxiliary definitions used throughout the appendix.

\begin{definition}[Sub-Gaussian Orlicz norm]\label{def:subgaussian}
For a random variable $Z$, the sub-Gaussian Orlicz norm is
\[
\|Z\|_{\psi_2}\;\coloneqq\;\inf\{t>0:\mathbb{E}[\exp(Z^2/t^2)]\le 2\},
\]
and $Z$ is called \emph{sub-Gaussian} when $\|Z\|_{\psi_2}<\infty$.
\end{definition}

\begin{definition}[Sparsity indices]\label{def:sparsity}
For the population targets of the calibrated losses, define the $\ell_0$-sparsity levels
\[
s_\alpha\coloneqq \|\boldsymbol{\alpha}^{1,*}\|_0,\qquad s_\gamma\coloneqq \|\boldsymbol{\gamma}_1^*\|_0,\qquad s\coloneqq s_\alpha+s_\gamma,
\]
and throughout the appendix we use $s$ (rather than the equivalent notation $s_{\mathrm{tot}}$) for the total sparsity. The same convention applies to the initial nuisance estimators under the two-level alignment of Assumption~\ref{asmp:rsc}(c).
\end{definition}

\begin{definition}[Target outcome regression]\label{def:mt1}
The target-site treated conditional mean is
\[
m_t^1(\boldsymbol{X})\;\coloneqq\;\mathbb{E}_t[Y\mid A=1,\boldsymbol{X}],
\]
so that under ignorability and overlap on the target site, $\mu^1=\mathbb{E}_t[m_t^1(\boldsymbol{X})]$.
\end{definition}

\begin{remark}[On $\phi(\boldsymbol{X})$]\label{rem:phi_convention}
Throughout, $\phi:\mathbb{R}^p\to\mathbb{R}^p$ denotes a basis expansion of the raw covariates $\boldsymbol{X}$ that includes an intercept coordinate and is treated as fixed. All high-dimensional and sparsity statements below apply to $\phi(\boldsymbol{X})$; if the user does not require an explicit expansion, $\phi$ can be taken as the identity, in which case $p=\dim(\boldsymbol{X})$.
\end{remark}

\begin{lemma}[Identification / double robustness for $\mu^1(\vartheta^*)$]\label{lem:ident_dr_mu}
Assume the target causal parameter is identified as $\mu^1=\E_t[Y(1)]$ (e.g., by consistency and ignorability/overlap on the target site), so that $\mu^1=\E_t[m_t^1(\boldsymbol{X})]$ where $m_t^1(\boldsymbol{X})\coloneqq \E_t[Y\mid A=1,\boldsymbol{X}]$.
Recall $\phi(\boldsymbol{X})$ includes an intercept.
Let $\vartheta^*=(\boldsymbol{\alpha}^{1,*},\boldsymbol{\gamma}_1^*)$ denote the population targets of the calibrated losses.
Then $\mu^1(\vartheta^*)=\mu^1$ holds under either of the following conditions.
\begin{enumerate}
    \item \textbf{Outcome-model correctness.} If $\psi(\phi(\boldsymbol{X});\boldsymbol{\alpha}^{1,*})=m_t^1(\boldsymbol{X})$ almost surely, then the calibrated outcome loss implies the population moment condition
    \(
    \E_s[w(\boldsymbol{X};\boldsymbol{\gamma}_1^*)\,\phi(\boldsymbol{X})\{Y-\psi(\phi(\boldsymbol{X});\boldsymbol{\alpha}^{1,*})\}]=\mathbf{0}
    \)
    (see Lemma~\ref{lem:ortho_foc}(ii) below). In particular, the intercept component yields
    \(
    \E_s[w(\boldsymbol{X};\boldsymbol{\gamma}_1^*)\{Y-\psi(\phi(\boldsymbol{X});\boldsymbol{\alpha}^{1,*})\}]=0
    \),
    and therefore $\mu^1(\vartheta^*)=\E_t[\psi(\phi(\boldsymbol{X});\boldsymbol{\alpha}^{1,*})]=\E_t[m_t^1(\boldsymbol{X})]=\mu^1$.
    \item \textbf{Weight-model correctness.} Suppose the weighting model is correct in the sense that for any square-integrable measurable function $h(\boldsymbol{X})$,
    \[
    \E_s[w(\boldsymbol{X};\boldsymbol{\gamma}_1^*)\,h(\boldsymbol{X})]=\E_t[h(\boldsymbol{X})],
    \]
    and, for the treated potential outcome, the usual covariate-shift/transportability condition holds for this pair $(t,s)$:
    \(
    \E_s[Y(1)\mid \boldsymbol{X}] = \E_t[Y(1)\mid \boldsymbol{X}]
    \)
    almost surely.
    Then for any outcome model $m(\boldsymbol{X})$,
    \[
    \E_s\!\left[w(\boldsymbol{X};\boldsymbol{\gamma}_1^*)\{Y(1)-m(\boldsymbol{X})\}\right]
    =
    \E_t\!\left[Y(1)-m(\boldsymbol{X})\right],
    \]
    and hence $\mu^1(\boldsymbol{\alpha},\boldsymbol{\gamma}_1^*)=\mu^1$ for any $\boldsymbol{\alpha}$, in particular $\mu^1(\vartheta^*)=\mu^1$ even if the outcome model is misspecified.
\end{enumerate}
\end{lemma}

Define the centered influence function at the population parameters $\vartheta^*$:
\begin{align*}
\varphi_{t,s}(D_i; \vartheta^*)
&:= \widetilde{\varphi}_{t,s}(D_i; \vartheta^*) - \mu^1(\vartheta^*) \\
&= \frac{I(R_i=t)}{\rho_t}\Big(\psi(\phi(\boldsymbol{X}_i); \boldsymbol{\alpha}^{1,*}) - \mu^1_t(\boldsymbol{\alpha}^{1,*})\Big) \\
&\quad+ \frac{I(R_i=s)}{\rho_s}\Big(
    w(\boldsymbol{X}_i;\boldsymbol{\gamma}_1^*)
    \{Y_i - \psi(\phi(\boldsymbol{X}_i); \boldsymbol{\alpha}^{1,*})\}
    - \delta^1_s(\vartheta^*)
\Big).
\end{align*}
where $\mu^1_t(\boldsymbol{\alpha}^{1,*}) = \E_t[\psi(\phi(\boldsymbol{X}); \boldsymbol{\alpha}^{1,*})]$ and $\delta^1_s(\vartheta^*) = \E_s[w(\boldsymbol{X};\boldsymbol{\gamma}_1^*) \{Y - \psi(\phi(\boldsymbol{X}); \boldsymbol{\alpha}^{1,*})\}]$.
More generally, for any $\vartheta$ we define $\varphi_{t,s}(D;\vartheta) \coloneqq \widetilde{\varphi}_{t,s}(D;\vartheta) - \mu^1(\vartheta)$.
Let $V^*_{t,s} := \Var\!\big(\varphi_{t,s}(D; \vartheta^*)\big)$.

\subsection{Assumptions}

The assumptions below are organized so that each condition has a clearly identified role: Assumption~\ref{asmp:data} controls the design distribution, Assumption~\ref{asmp:overlap} controls the weight tails together with the outcome smoothness and moments, and Assumption~\ref{asmp:rsc} controls the nuisance estimation pipeline. Definitions of Orlicz norms, sparsity, and the target outcome regression are collected in Section~\ref{sec:notation_defs}.

\begin{assumption}[Design distribution]\label{asmp:data}
The coordinates of $\phi(\boldsymbol{X})$ (which includes an intercept) are sub-Gaussian in the sense of Definition~\ref{def:subgaussian}: there exists $K_\phi<\infty$ such that
\[
\max_{1\le j\le p}\|\phi_j(\boldsymbol{X})\|_{\psi_2}\le K_\phi\qquad\text{under both }R=t\text{ and }R=s.
\]
\end{assumption}

\begin{assumption}[Overlap, smoothness, and moments]\label{asmp:overlap}
\begin{enumerate}[label={(\alph*)},leftmargin=*]
    \item (Bounded-weight overlap) There exists a constant $\bar{W}<\infty$ such that the ideal calibration weight at the population target is essentially bounded:
    \[
    w(\boldsymbol{X};\boldsymbol{\gamma}_1^*)\le \bar{W}\quad\text{almost surely under }R=s,
    \]
    equivalently, $\phi(\boldsymbol{X})^\top\boldsymbol{\gamma}_1^* \ge -\log \bar{W}$ a.s. Under this condition, numerical truncation in the calibrated losses is asymptotically inactive at the target.

    \item (GLM smoothness) The mean function $\psi:\mathbb{R}\to\mathbb{R}$ is continuously differentiable with $\sup_u|\psi'(u)|\le L_\psi<\infty$, and $\psi'$ is Lipschitz with constant $L_{\psi'}<\infty$.

    \item (Outcome moments) There exists $\delta>0$ such that $\mathbb{E}[|Y|^{2+\delta}\mid R=r]<\infty$ for $r\in\{t,s\}$. A sufficient condition, used in the concentration steps below, is that $Y\mid(A=1,R=r)$ is sub-Gaussian.

    \item (Local IF moments) With the same $\delta$ as in (c), there exists a neighborhood $\mathcal{N}(\vartheta^*)$ of the population target such that the uncentered influence-function contribution $\widetilde{\varphi}_{t,s}$ defined in the Setup has a uniformly bounded $(2+\delta)$-th moment,
    \[
    \sup_{\vartheta\in\mathcal{N}(\vartheta^*)}\mathbb{E}\!\left[|\widetilde{\varphi}_{t,s}(D;\vartheta)|^{2+\delta}\right]<\infty.
    \]
    Under (a), (b), (c) and Assumption~\ref{asmp:data}, this condition is automatic: the bounded weight and sub-Gaussian $\phi(\boldsymbol{X})$ imply finite moments of $w(\boldsymbol{X};\boldsymbol{\gamma})\,\psi(\phi(\boldsymbol{X})^\top\boldsymbol{\alpha})$ of any order in a neighborhood of $\vartheta^*$.

    \item (Inactive truncation) With $M_{\tau,N}$ as in the definition of $\mathcal{T}(\cdot)$, for $r\in\{t,s\}$,
    \[
    P\!\left(\bigl|\phi(\boldsymbol{X})^\top\boldsymbol{\alpha}^{1,*}\bigr| > M_{\tau,N}\,\Big|\, R=r\right)=o(N^{-1}),\qquad
    P\!\left(\bigl|\phi(\boldsymbol{X})^\top\boldsymbol{\gamma}_1^*\bigr| > M_{\tau,N}\,\Big|\, R=s\right)=o(N^{-1}).
    \]
    This is stronger than strictly necessary but keeps the exposition clean; weakening it to $o(1)$ is possible at the cost of tracking a clipping-induced bias term of order $o_p(N^{-1/2})$.
\end{enumerate}
\end{assumption}

\begin{assumption}[Nuisance estimation pipeline]\label{asmp:rsc}
\begin{enumerate}[label={(\alph*)},leftmargin=*]
    \item (Restricted strong convexity) The population versions of the calibrated loss functions (Equations~\eqref{eq:gamma_calibrated_loss} and \eqref{eq:alpha_calibrated_loss} in the main manuscript) are convex and satisfy Restricted Strong Convexity (RSC) around $\vartheta^*$ in the relevant sparse norms.
    \item (Cross-fitting) The estimators $\widehat{\vartheta}^{(-k)}$ used for evaluation on fold $k$ are trained solely on $\mathcal{D}^{(-k)}$, which is independent of $\mathcal{D}_k$.
    \item (Two-level limit alignment) The probability limits $(\boldsymbol{\alpha}^{1,*}_{\mathrm{init}},\boldsymbol{\gamma}^{*}_{1,\mathrm{init}})$ of the initial plug-in nuisance estimators coincide with the population targets of the calibrated step, i.e., $\boldsymbol{\alpha}^{1,*}_{\mathrm{init}}=\boldsymbol{\alpha}^{1,*}$ and $\boldsymbol{\gamma}^{*}_{1,\mathrm{init}}=\boldsymbol{\gamma}^{*}_{1}$.
    \item (Foldwise score-centering; optional) The conditional foldwise score-centering condition of the \emph{Special} case in Lemma~\ref{lem:nuisance_rates} holds for the calibrated losses used to estimate $\boldsymbol{\alpha}^1$ and $\boldsymbol{\gamma}_1$; this is the condition exploited by \citet{hou2025efficient} to eliminate the plug-in score term. When (d) is not verified, the \emph{General} rate bound in Lemma~\ref{lem:nuisance_rates} applies instead.
\end{enumerate}
\end{assumption}

\begin{remark}[On the alignment condition]\label{rem:alignment}
Assumption \ref{asmp:rsc}(c) is a technical device to avoid tracking multiple population limits induced by the nested/coupled optimization. It is satisfied, for example, when (i) the population calibrated losses have unique minimizers $(\boldsymbol{\alpha}^{1,*},\boldsymbol{\gamma}_1^*)$ (by strict convexity in the relevant restricted set), and (ii) the mapping from an initial plug-in value to the population minimizer of the next calibrated update is continuous at $(\boldsymbol{\alpha}^{1,*},\boldsymbol{\gamma}_1^*)$, so that any consistent initialization converges to the same fixed point. In practice, the same calibrated estimating equations are solved throughout the nuisance-estimation pipeline, making the alignment natural under uniqueness.
\end{remark}

\subsection{Oracle Rates for the Calibrated Nuisance Estimators}\label{sec:oracle}

Before stating the asymptotic normality of the cross-fitted estimator, we record the high-dimensional oracle rates of the calibrated nuisance estimators. The result is a direct specialization of the cross-fitted $\ell_1$-penalized nuisance rate bound of Lemma~\ref{lem:nuisance_rates} (which in turn adapts Lemma~A21 of \citet{hou2025efficient} and the RSC framework of \citet{negahban2012unified,tan2020model,buhlmann2011statistics}) to the calibrated outcome and calibration-weight losses of Equations~\eqref{eq:alpha_calibrated_loss} and \eqref{eq:gamma_calibrated_loss}. Separating this oracle step from the subsequent asymptotic-normality argument keeps the role of the nuisance rates logically distinct from the debiasing argument and facilitates reuse across the pairwise and aggregated statements.

\begin{proposition}[Oracle rates of calibrated nuisance estimators]\label{prop:oracle_rates}
Under Assumptions~\ref{asmp:data}, \ref{asmp:overlap}, and \ref{asmp:rsc}(a)--(c), with regularization parameters $\lambda_\alpha\asymp\sqrt{\log(p)/N}$ and $\lambda_\gamma\asymp\sqrt{\log(p)/N}$, the calibrated cross-fitted $\ell_1$-penalized nuisance estimators satisfy
\[
\|\widehat{\boldsymbol{\alpha}}_{t,s}^{1,(-k)}-\boldsymbol{\alpha}^{1,*}\|_2 = O_p\!\left(\sqrt{\frac{s_\alpha\log(p)}{N}}\right),
\quad
\|\widehat{\boldsymbol{\alpha}}_{t,s}^{1,(-k)}-\boldsymbol{\alpha}^{1,*}\|_1 = O_p\!\left(s_\alpha\sqrt{\frac{\log(p)}{N}}\right),
\]
and analogously for $\widehat{\boldsymbol{\gamma}}_{s,1}^{(-k)}$ with $s_\alpha$ replaced by $s_\gamma$. If Assumption~\ref{asmp:rsc}(d) additionally holds, the foldwise score-centering is available and the \emph{Special} rate bound in Lemma~\ref{lem:nuisance_rates} applies without inflation from the plug-in score term.
\end{proposition}

\begin{proof}
The proposition is obtained by applying Lemma~\ref{lem:nuisance_rates} twice, once to the calibrated outcome loss (Equation~\eqref{eq:alpha_calibrated_loss} in the main manuscript) and once to the calibrated calibration-weight loss (Equation~\eqref{eq:gamma_calibrated_loss}). We verify each condition of Lemma~\ref{lem:nuisance_rates} and then derive the stated rates.

\textbf{Step 1 (Outcome loss as an instance of Lemma~\ref{lem:nuisance_rates}).}
Identify the generic estimation target of Lemma~\ref{lem:nuisance_rates} with $\boldsymbol{\beta}\equiv\boldsymbol{\alpha}^1$, the plug-in nuisance with $\boldsymbol{\gamma}\equiv\boldsymbol{\gamma}_1$, and the per-observation loss with
\[
\ell(D;\boldsymbol{\alpha}^1,\boldsymbol{\gamma}_1)=\frac{I(R=s)}{\rho_s}\frac{I(A=1)}{e^{\mathcal{T}(\phi(\boldsymbol{X})^\top\boldsymbol{\gamma}_1)}}\Bigl[-Y\,\phi(\boldsymbol{X})^\top\boldsymbol{\alpha}^1+\Psi(\phi(\boldsymbol{X})^\top\boldsymbol{\alpha}^1)\Bigr].
\]
We check the conditions (a)--(d) of Lemma~\ref{lem:nuisance_rates}:
\begin{enumerate}[label={(\alph*)},leftmargin=*]
\item Sub-Gaussian design with bounded $\psi_2$-norm is supplied by Assumption~\ref{asmp:data}.
\item Response sub-Gaussianity is Assumption~\ref{asmp:overlap}(c); the GLM mean function $\psi$ and its derivative $\psi'$ are uniformly bounded on a neighborhood of $\boldsymbol{\alpha}^{1,*}$ by Assumption~\ref{asmp:overlap}(b), using that $\phi(\boldsymbol{X})^\top\boldsymbol{\alpha}^{1,*}$ has a sub-Gaussian distribution by Assumption~\ref{asmp:data}.
\item Identifiability of the population Gram matrix $\mathbb{E}[\phi(\boldsymbol{X})\phi(\boldsymbol{X})^\top\mid R=s]$ follows from the RSC of the outcome calibrated loss (Assumption~\ref{asmp:rsc}(a)): RSC with curvature $\kappa_\alpha>0$ implies that the restricted minimum eigenvalue of the population Hessian at $\boldsymbol{\alpha}^{1,*}$ is bounded below, which, combined with bounded weights (Assumption~\ref{asmp:overlap}(a)), yields a strictly positive minimum eigenvalue of the unrestricted Gram matrix on the support of $\boldsymbol{\alpha}^{1,*}$ and in its local extensions.
\item Weight regularity: the calibration weight $w(\boldsymbol{X};\boldsymbol{\gamma})$ is locally Lipschitz in $\boldsymbol{\gamma}$ because $u\mapsto e^{-u}$ is smooth, and its moments are finite under (a) of Assumption~\ref{asmp:overlap}.
\end{enumerate}
Cross-fitting independence is given by Assumption~\ref{asmp:rsc}(b), and the alignment condition (c) guarantees that the initial plug-in $\widehat{\boldsymbol{\gamma}}_1^{(-k)}$ converges to the same population target $\boldsymbol{\gamma}_1^*$ used in the analysis.
With $\lambda_\alpha\asymp\sqrt{\log(p)/N}$, Lemma~\ref{lem:nuisance_rates} yields the stated $\ell_2$ and $\ell_1$ rates for $\widehat{\boldsymbol{\alpha}}_{t,s}^{1,(-k)}-\boldsymbol{\alpha}^{1,*}$ with the sparsity index $s_\alpha$ of Definition~\ref{def:sparsity}.

\textbf{Step 2 (Calibration-weight loss).}
Identify the generic target with $\boldsymbol{\beta}\equiv\boldsymbol{\gamma}_1$, the plug-in nuisance with $\boldsymbol{\gamma}\equiv\boldsymbol{\alpha}^1$, and the per-observation loss with
\[
\ell(D;\boldsymbol{\gamma}_1,\boldsymbol{\alpha}^1)=\frac{I(R=t)}{\rho_t}\psi'\!\left(\mathcal{T}(\phi(\boldsymbol{X})^\top\boldsymbol{\alpha}^1)\right)\phi(\boldsymbol{X})^\top\boldsymbol{\gamma}_1+\frac{I(R=s)}{\rho_s}\frac{I(A=1)\,\psi'\!\left(\mathcal{T}(\phi(\boldsymbol{X})^\top\boldsymbol{\alpha}^1)\right)}{e^{\phi(\boldsymbol{X})^\top\boldsymbol{\gamma}_1}}.
\]
The loss is strictly convex in $\boldsymbol{\gamma}_1$ because $u\mapsto e^{-u}$ is strictly convex, and its Hessian at $\boldsymbol{\gamma}_1^*$ equals $\mathbb{E}_s[w(\boldsymbol{X};\boldsymbol{\gamma}_1^*)\psi'(\phi(\boldsymbol{X})^\top\boldsymbol{\alpha}^{1,*})\phi(\boldsymbol{X})\phi(\boldsymbol{X})^\top]/\rho_s$, which is positive definite on the support of $\boldsymbol{\gamma}_1^*$ by Assumption~\ref{asmp:rsc}(a). Conditions (a)--(d) of Lemma~\ref{lem:nuisance_rates} are verified exactly as in Step~1 (with $\boldsymbol{\alpha}^1$ and $\boldsymbol{\gamma}_1$ swapped), using that $\sup_u|\psi'(u)|\le L_\psi$ (Assumption~\ref{asmp:overlap}(b)) and the bounded-weight condition (Assumption~\ref{asmp:overlap}(a)). With $\lambda_\gamma\asymp\sqrt{\log(p)/N}$, Lemma~\ref{lem:nuisance_rates} delivers the stated rates for $\widehat{\boldsymbol{\gamma}}_{s,1}^{(-k)}-\boldsymbol{\gamma}_1^*$ with sparsity index $s_\gamma$.

\textbf{Step 3 (Foldwise score-centering).}
If Assumption~\ref{asmp:rsc}(d) additionally holds, the foldwise conditional score-centering condition required by the \emph{Special} branch of Lemma~\ref{lem:nuisance_rates} is satisfied for both calibrated losses, and the corresponding rate bound applies without the plug-in score term; otherwise, the \emph{General} branch of the lemma applies, giving the same stated rates up to constants.
\end{proof}

The asymptotic-normality statements in Sections~\ref{sec:main_normality} and \ref{sec:agg_proofs} will invoke Proposition~\ref{prop:oracle_rates} exclusively through the product rate
\[
\|\widehat{\boldsymbol{\alpha}}_{t,s}^{1,(-k)}-\boldsymbol{\alpha}^{1,*}\|_2\cdot\|\widehat{\boldsymbol{\gamma}}_{s,1}^{(-k)}-\boldsymbol{\gamma}_1^*\|_2 = O_p\!\left(\frac{\sqrt{s_\alpha s_\gamma}\log(p)}{N}\right),
\]
which, combined with the cross-calibration property of Lemma~\ref{lem:ortho_foc}, drives the bias remainder $R_{2,k}$ to $o_p(N^{-1/2})$.

\subsection{Main Theorem and Proof}\label{sec:main_normality}

\begin{theorem}[Asymptotic Normality of $\widehat{\mu}^{1,\mathrm{CF}}_{t,s}$]\label{thm:main}
Suppose Assumptions \ref{asmp:data}, \ref{asmp:overlap}, and \ref{asmp:rsc}(a)--\ref{asmp:rsc}(c) hold. (Assumption \ref{asmp:rsc}(d) is an optional strengthening used only to invoke the \textbf{Special} rate in Lemma \ref{lem:nuisance_rates}; the theorem also holds under the \textbf{General} rate condition there.) Assume $\vartheta^*$ is the unique population target pair induced by the calibrated losses under the alignment condition in Assumption \ref{asmp:rsc}(c), implying Lemma \ref{lem:ortho_foc} holds. If further, at least one nuisance model (outcome or weighting) is correctly specified in the sense of Lemma \ref{lem:ident_dr_mu} (so $\mu^1(\vartheta^*)=\mu^1$), and the total sparsity $s=s_\alpha+s_\gamma$ of Definition~\ref{def:sparsity} (applied to both the initial and the calibrated nuisance targets under the alignment condition of Assumption~\ref{asmp:rsc}(c)) satisfies
\[
\frac{s^2\, \log^2(p)}{N}\;\longrightarrow\;0,
\]
then the cross-fitted estimator is asymptotically normal:
\[
\sqrt{N}\,\big(\widehat{\mu}^{1,\mathrm{CF}}_{t,s} - \mu^1\big)
\;\xrightarrow{d}\; \mathcal{N}\!\big(0,\,V^*_{t,s}\big).
\]
Moreover, the finite-sample KKT perturbation from the $\ell_1$ penalties is absorbed in the remainder term (via Lemma~\ref{lem:nuisance_rates}) and is $o_p(N^{-1/2})$ under the stated sparsity/rate condition.
\end{theorem}

\begin{proof}
Suppose in addition that at least one nuisance model is correctly specified in the sense of Lemma \ref{lem:ident_dr_mu}, so that $\mu^1(\vartheta^*)=\mu^1$. We may then center the expansion at $\mu^1$.
The total estimation error is decomposed as:
Write the decomposition
\[
\begin{aligned}
\widehat{\mu}^{1,\mathrm{CF}}_{t,s} - \mu^1
&= \underbrace{\left(\frac{1}{N}\sum_{i=1}^{N} \varphi_{t,s}(D_i; \vartheta^*)\right)}_{\displaystyle \mathsf{T}_{\mathrm{I}}}
+ \underbrace{\left(\widehat{\mu}^{1,\mathrm{CF}}_{t,s} - \frac{1}{N}\sum_{i=1}^{N} \varphi_{t,s}(D_i; \vartheta^*) - \mu^1\right)}_{\displaystyle \mathsf{T}_{\mathrm{II}}},
\end{aligned}
\]
where $\mathsf{T}_{\mathrm{I}}$ is the average of centered influence functions at the population target and $\mathsf{T}_{\mathrm{II}}$ is the remainder; $\mathbb{E}[\varphi_{t,s}(D;\vartheta^*)]=0$ by construction and $\mu^1(\vartheta^*)=\mu^1_t(\boldsymbol{\alpha}^{1,*})+\delta^1_s(\vartheta^*)$.
By Lemma~\ref{lem:clt},
\[
\frac{1}{\sqrt{N}}\sum_{i=1}^{N}\varphi_{t,s}(D_i; \vartheta^*) \xrightarrow{d} \mathcal{N}(0, V^*_{t,s}),
\]
so $\sqrt{N}\,\mathsf{T}_{\mathrm{I}}\xrightarrow{d}\mathcal{N}(0,V^*_{t,s})$.
The remainder is an average over the folds, $\mathsf{T}_{\mathrm{II}}=\frac{1}{K_f}\sum_{k=1}^{K_f}(R_{1,k}+R_{2,k})$, with components analyzed below.
$R_{1,k}$ is the conditional empirical process fluctuation of the \emph{difference} score:
\[
R_{1,k} = \left(\tildeE_k - \E[\cdot \mid \mathcal{D}^{(-k)}]\right)\left[
\widetilde{\varphi}_{t,s}(D; \widehat{\vartheta}^{(-k)}) - \widetilde{\varphi}_{t,s}(D; \vartheta^*)
\right].
\]
By cross-fitting, conditional on $\mathcal{D}^{(-k)}$, the summands are independent (with a two-sample/triangular-array structure) on fold $k$ with mean zero. Using Lemma \ref{lem:bias_bound} (Part 2),
\[
\E\!\left[\left(\widetilde{\varphi}_{t,s}(D; \widehat{\vartheta}^{(-k)}) - \widetilde{\varphi}_{t,s}(D; \vartheta^*)\right)^2 \,\middle|\, \mathcal{D}^{(-k)}\right]
\le C\!\left(\|\widehat{\boldsymbol{\alpha}}_{t,s}^{1,(-k)}-\boldsymbol{\alpha}^{1,*}\|_2^2
\!+\!
\|\widehat{\boldsymbol{\gamma}}_{s,1}^{(-k)}-\boldsymbol{\gamma}_{1}^{*}\|_2^2\right),
\]
and therefore the conditional variance of $R_{1,k}$ obeys
\[
\Var(R_{1,k}\mid \mathcal{D}^{(-k)})
=
\Var\!\left(
\tildeE_k\!\left[\Delta_k(D)\right]\middle|\mathcal{D}^{(-k)}
\right)
\le
\frac{1}{N_k}\E\!\left[\Delta_k(D)^2\mid \mathcal{D}^{(-k)}\right],
\]
where $\Delta_k(D)\coloneqq \widetilde{\varphi}_{t,s}(D; \widehat{\vartheta}^{(-k)}) - \widetilde{\varphi}_{t,s}(D; \vartheta^*)$.
Here we used that, conditional on $\mathcal{D}^{(-k)}$, the $\{\Delta_k(D_i): i\in\mathcal{D}_k\}$ are independent, so
\[
\begin{aligned}
\Var\!\left(\tildeE_k[\Delta_k(D)]\mid\mathcal{D}^{(-k)}\right)
&=
\Var\!\left(\tildeE_k[\Delta_k(D)]-\E[\Delta_k(D)\mid\mathcal{D}^{(-k)}]\mid\mathcal{D}^{(-k)}\right) \\
&\le
\frac{1}{N_k}\E\!\left[(\Delta_k(D)-\E[\Delta_k(D)\mid\mathcal{D}^{(-k)}])^2\mid\mathcal{D}^{(-k)}\right] \\
&\le
\frac{1}{N_k}\E\!\left[\Delta_k(D)^2\mid\mathcal{D}^{(-k)}\right].
\end{aligned}
\]
Since $N_k\asymp N$ for fixed $K_f$, Chebyshev's inequality yields
\[
|R_{1,k}| = O_p\!\left(\frac{
\|\widehat{\boldsymbol{\alpha}}_{t,s}^{1,(-k)}-\boldsymbol{\alpha}^{1,*}\|_2
\!+\!
\|\widehat{\boldsymbol{\gamma}}_{s,1}^{(-k)}-\boldsymbol{\gamma}_{1}^{*}\|_2
}{\sqrt{N}}\right)
= O_p\!\left(\frac{\sqrt{s\log(p)}}{N}\right)
= o_p(N^{-1/2}).
\]

$R_{2,k}$ is the first-order bias remainder, conditional on $\mathcal{D}^{(-k)}$:
\[
R_{2,k} = \E[\widetilde{\varphi}_{t,s}(D; \widehat{\vartheta}^{(-k)}) - \widetilde{\varphi}_{t,s}(D; \vartheta^*) \mid \mathcal{D}^{(-k)}]
\]
By Lemma \ref{lem:bias_bound} (Part 1), which relies on the orthogonality conditions derived in Lemma \ref{lem:ortho_foc}, this bias is bounded by a second-order term:
\[
|R_{2,k}| \le C \cdot \left(
\lVert \widehat{\boldsymbol{\alpha}}_{t,s}^{1,(-k)} - \boldsymbol{\alpha}^{1,*} \rVert_2^2
+
\lVert \widehat{\boldsymbol{\gamma}}_{s,1}^{(-k)} - \boldsymbol{\gamma}_1^* \rVert_2^2
+
\lVert \widehat{\boldsymbol{\alpha}}_{t,s}^{1,(-k)} - \boldsymbol{\alpha}^{1,*} \rVert_2
\lVert \widehat{\boldsymbol{\gamma}}_{s,1}^{(-k)} - \boldsymbol{\gamma}_1^* \rVert_2
\right)
\]
Using the nuisance rates from Lemma \ref{lem:nuisance_rates}:
\[
|R_{2,k}| = O_p\left(\frac{s_\alpha \log(p)}{N}\right) + O_p\left(\frac{s_\gamma \log(p)}{N}\right) + O_p\left( \frac{\sqrt{s_\alpha s_\gamma} \log(p)}{N} \right)
= O_p\left( \frac{s \log(p)}{N} \right).
\]
Averaging over $k$ gives the total bias remainder:
\[
R_2 = \frac{1}{K_f}\sum_{k=1}^{K_f} R_{2,k} = O_p\!\left( \frac{s \log(p)}{N} \right).
\]
We require $\sqrt{N}\cdot R_2 \to 0$, i.e., $O_p(s \log(p) / \sqrt{N}) \to 0$, which is implied by the theorem's condition $s^2 \log^2(p)/N \to 0$.
Thus the total remainder $\mathsf{T}_{\mathrm{II}}=R_1+R_2$ satisfies $\sqrt{N}\,\mathsf{T}_{\mathrm{II}}=o_p(1)$, and hence
\[
\sqrt{N}\,\big(\widehat{\mu}^{1,\mathrm{CF}}_{t,s} - \mu^1\big)
=
\frac{1}{\sqrt{N}}\sum_{i=1}^{N}\varphi_{t,s}(D_i;\vartheta^*) + o_p(1).
\]
Together with Lemma \ref{lem:clt}, this yields $\sqrt{N}\,(\widehat{\mu}^{1,\mathrm{CF}}_{t,s} - \mu^1)\xrightarrow{d}\mathcal{N}(0,V^*_{t,s})$ by Slutsky's theorem.
\end{proof}

\begin{lemma}[Consistent variance estimation]\label{lem:var_est}
Assume Assumptions~\ref{asmp:data}--\ref{asmp:rsc} and that the population asymptotic variance is bounded away from zero, $V^*_{t,s}\ge \underline{v}>0$. Then the plug-in estimator based on the cross-fitted empirical influence function is consistent:
\[
\widehat{V}_{t,s} = \frac{1}{K_f} \sum_{k=1}^{K_f} \tildeE_k \left[ \left( \widetilde{\varphi}_{t,s}(D; \widehat{\vartheta}^{(-k)}) - \widehat{\mu}^{1,\mathrm{CF}}_{t,s} \right)^2 \right] \xrightarrow{p} V^*_{t,s}.
\]
\end{lemma}

\begin{proof}
The argument combines three ingredients: the nuisance consistency rates of Proposition~\ref{prop:oracle_rates} (equivalently, $\|\widehat{\vartheta}^{(-k)}-\vartheta^*\|_2=o_p(1)$), the local $L_2$-Lipschitz continuity of $\vartheta\mapsto\widetilde{\varphi}_{t,s}(D;\vartheta)$ (there exists $C<\infty$ such that, for $\vartheta_1,\vartheta_2$ in a neighborhood of $\vartheta^*$, $\mathbb{E}[(\widetilde{\varphi}_{t,s}(D;\vartheta_1)-\widetilde{\varphi}_{t,s}(D;\vartheta_2))^2]\le C\|\vartheta_1-\vartheta_2\|_2^2$, which follows from Assumptions~\ref{asmp:overlap}(a)--(b) as in Lemma~\ref{lem:bias_bound}(Part~2)), and uniform integrability of $\{\widetilde{\varphi}_{t,s}(D;\widehat{\vartheta}^{(-k)})^2\}_{k,N}$ (supplied by the $(2+\delta)$-moment bound in Assumption~\ref{asmp:overlap}(d)).

\textbf{Step 1 (Reduce to within-fold centering).}
	For each fold $k$, define the foldwise score
	\[
	Z_{i}^{(k)} \coloneqq \widetilde{\varphi}_{t,s}(D_i;\widehat{\vartheta}^{(-k)}),
	\qquad i\in\mathcal{D}_k,
	\]
		and the foldwise mean $\widehat{\mu}_k\coloneqq \tildeE_k[Z^{(k)}]$.
		Since $\widehat{\mu}^{1,\mathrm{CF}}_{t,s}=\frac{1}{K_f}\sum_{k=1}^{K_f}\widehat{\mu}_k$, we have the identity
		\[
		\tildeE_k\!\left[(Z^{(k)}-\widehat{\mu}^{1,\mathrm{CF}}_{t,s})^2\right]
		=
		\tildeE_k\!\left[(Z^{(k)}-\widehat{\mu}_k)^2\right]
		+
		\left(\widehat{\mu}_k-\widehat{\mu}^{1,\mathrm{CF}}_{t,s}\right)^2,
		\]
		where the cross term vanishes because $\tildeE_k[Z^{(k)}-\widehat{\mu}_k]=0$ by definition of $\widehat{\mu}_k$.
		Therefore,
		\[
		\widehat{V}_{t,s}
		=
		\frac{1}{K_f}\sum_{k=1}^{K_f}\tildeE_k\!\left[(Z^{(k)}-\widehat{\mu}_k)^2\right]
		+
		\frac{1}{K_f}\sum_{k=1}^{K_f}\left(\widehat{\mu}_k-\widehat{\mu}^{1,\mathrm{CF}}_{t,s}\right)^2.
		\]
	By Lemma \ref{lem:nuisance_rates}, 
    $\widehat{\vartheta}^{(-k)}=\vartheta^*+o_p(1)$, and by the conditional LLN below we have $\widehat{\mu}_k=\mu^1(\vartheta^*)+o_p(1)$ for each fixed $k$. Since $K_f$ is fixed, this implies
	\(
	\max_{1\le k\le K_f}\left|\widehat{\mu}_k-\widehat{\mu}^{1,\mathrm{CF}}_{t,s}\right| = o_p(1)
	\),
	and hence the second term is $o_p(1)$. Thus it suffices to show
	\[
	\frac{1}{K_f}\sum_{k=1}^{K_f}\tildeE_k\!\left[(Z^{(k)}-\widehat{\mu}_k)^2\right]\xrightarrow{p}V^*_{t,s}.
	\]

\textbf{Step 2 (Conditional LLN for first and second moments).}
Fix a fold $k$. Conditional on the training sample $\mathcal{D}^{(-k)}$, the nuisance estimate $\widehat{\vartheta}^{(-k)}$ is fixed and independent of $\mathcal{D}_k$.
Therefore, conditional on $\mathcal{D}^{(-k)}$, the random variables $\{Z_i^{(k)}:i\in\mathcal{D}_k\}$ are independent with finite $(2+\delta)$ moments uniformly in $N$ by Assumption \ref{asmp:overlap}(b).
By the (conditional) weak law of large numbers,
\[
\widehat{\mu}_k - \E\!\left[Z^{(k)}\mid \mathcal{D}^{(-k)}\right] = o_p(1),
\qquad
\tildeE_k\!\left[(Z^{(k)})^2\right] - \E\!\left[(Z^{(k)})^2\mid \mathcal{D}^{(-k)}\right] = o_p(1).
\]
Since $\tildeE_k[(Z^{(k)}-\widehat{\mu}_k)^2]=\tildeE_k[(Z^{(k)})^2]-\widehat{\mu}_k^{\,2}$, combining the two displays yields
\[
\tildeE_k\!\left[(Z^{(k)}-\widehat{\mu}_k)^2\right]
-
\left\{
\E\!\left[(Z^{(k)})^2\mid \mathcal{D}^{(-k)}\right]
-
\E\!\left[Z^{(k)}\mid \mathcal{D}^{(-k)}\right]^2
\right\}
=o_p(1).
\]

\textbf{Step 3 (Plug-in stability: conditional expectations converge to population targets).}
It remains to show that the bracketed conditional variance converges in probability to $V^*_{t,s}$.
First, by nuisance consistency (Ingredient 1) and the $L_2$-Lipschitz property (Ingredient 2),
\[
\E\!\left[\left(Z^{(k)}-\widetilde{\varphi}_{t,s}(D;\vartheta^*)\right)^2\right]
=
\E\!\left[\left(\widetilde{\varphi}_{t,s}(D;\widehat{\vartheta}^{(-k)})-\widetilde{\varphi}_{t,s}(D;\vartheta^*)\right)^2\right]
\le C\|\widehat{\vartheta}^{(-k)}-\vartheta^*\|_2^2
\xrightarrow{p}0.
\]
Next we establish convergence of the \emph{conditional} first and second moments.
Let
\[
m_k \coloneqq \E\!\left[Z^{(k)}\mid \mathcal{D}^{(-k)}\right],
\qquad
q_k \coloneqq \E\!\left[(Z^{(k)})^2\mid \mathcal{D}^{(-k)}\right],
\qquad
\mu^* \coloneqq \mu^1(\vartheta^*)=\E[\widetilde{\varphi}_{t,s}(D;\vartheta^*)].
\]
Then
\[
m_k-\mu^*
=
\E\!\left[\widetilde{\varphi}_{t,s}(D;\widehat{\vartheta}^{(-k)})-\widetilde{\varphi}_{t,s}(D;\vartheta^*)\;\middle|\;\mathcal{D}^{(-k)}\right],
\]
so by conditional Jensen,
\[
(m_k-\mu^*)^2
\le
\E\!\left[\left(\widetilde{\varphi}_{t,s}(D;\widehat{\vartheta}^{(-k)})-\widetilde{\varphi}_{t,s}(D;\vartheta^*)\right)^2\;\middle|\;\mathcal{D}^{(-k)}\right].
\]
Taking expectations and using the $L_2$-Lipschitz bound in Ingredient 2 gives $\E[(m_k-\mu^*)^2]\to 0$, hence $m_k\to_p \mu^*$.
Similarly, write
\[
q_k-\E[\widetilde{\varphi}_{t,s}(D;\vartheta^*)^2]
=
\E\!\left[(Z^{(k)})^2-\widetilde{\varphi}_{t,s}(D;\vartheta^*)^2\;\middle|\;\mathcal{D}^{(-k)}\right].
\]
Using $(a^2-b^2)=(a-b)(a+b)$ and conditional Cauchy--Schwarz,
\[
\begin{aligned}
|q_k-\E[\widetilde{\varphi}_{t,s}(D;\vartheta^*)^2]|
&=
\left|
\E\!\left[(Z^{(k)})^2-\widetilde{\varphi}_{t,s}(D;\vartheta^*)^2\;\middle|\;\mathcal{D}^{(-k)}\right]
\right| \\
&\le
\E\!\left[\left|Z^{(k)}-\widetilde{\varphi}_{t,s}(D;\vartheta^*)\right|\cdot\left|Z^{(k)}+\widetilde{\varphi}_{t,s}(D;\vartheta^*)\right|\;\middle|\;\mathcal{D}^{(-k)}\right] \\
&\le
\left(
\E\!\left[\left(Z^{(k)}-\widetilde{\varphi}_{t,s}(D;\vartheta^*)\right)^2\;\middle|\;\mathcal{D}^{(-k)}\right]
\right)^{1/2}
\left(
\E\!\left[\left(Z^{(k)}+\widetilde{\varphi}_{t,s}(D;\vartheta^*)\right)^2\;\middle|\;\mathcal{D}^{(-k)}\right]
\right)^{1/2}.
\end{aligned}
\]
The first factor converges to 0 in probability by the $L_2$-Lipschitz bound in Ingredient 2 and nuisance consistency (Ingredient 1).
The second factor is $O_p(1)$ by the $(2+\delta)$ moment assumption in Assumption \ref{asmp:overlap}(b) (uniform integrability).
Therefore $q_k\to_p \E[\widetilde{\varphi}_{t,s}(D;\vartheta^*)^2]$.
Consequently, since $m_k\to_p \mu^*$,
\[
\E\!\left[(Z^{(k)})^2\mid \mathcal{D}^{(-k)}\right]
-
\E\!\left[Z^{(k)}\mid \mathcal{D}^{(-k)}\right]^2
\xrightarrow{p}
\E\!\left[\widetilde{\varphi}_{t,s}(D;\vartheta^*)^2\right]-\mu^1(\vartheta^*)^2
=
\Var\!\left(\widetilde{\varphi}_{t,s}(D;\vartheta^*)\right)
=V^*_{t,s}.
\]

\textbf{Step 4 (Average over folds).}
Since $K_f$ is fixed, combining Step 2 and Step 3 and averaging over $k$ yields $\widehat{V}_{t,s}\xrightarrow{p}V^*_{t,s}$.

\textbf{Extension to the aggregated estimator.}
For the weighted aggregated estimator with inner $M$-fold cross-fitted weights $\widehat{\boldsymbol{\eta}}^{(-k)}$, we define per-observation \emph{uncentered} scaled influence functions that serve as the unified estimating equation for both the point estimate and the variance.

\textbf{Unified estimating equation.}
For each observation $i$ evaluated on fold $\mathcal{D}_{k(i)}$, define the uncentered IF
\[
\widehat{\Phi}_i \coloneqq
\begin{cases}
\displaystyle\frac{N_{\mathrm{all}}}{N_t}\biggl[(1-\textstyle\sum_j \widehat{\eta}_j^{(-k(i))})\,\widetilde{Y}_{\mathrm{ot},i} + \textstyle\sum_j \widehat{\eta}_j^{(-k(i))}\,\psi(\phi(\bX_i);\widehat{\balpha}_{t,s_j}^{1,(-k(i))})\biggr], & i\in\mathcal{D}_t,\\[8pt]
\displaystyle\frac{N_{\mathrm{all}}}{N_{s_j}}\,\widehat{\eta}_j^{(-k(i))}\,w(\bX_i;\widehat{\bgamma}_{s_j,1}^{(-k(i))})\bigl(Y_i - \psi(\phi(\bX_i);\widehat{\balpha}_{t,s_j}^{1,(-k(i))})\bigr), & i\in\mathcal{D}_{s_j},
\end{cases}
\]
Here $\widetilde{Y}_{\mathrm{ot},i}$ is the \emph{uncentered} AIPW pseudo-outcome for the target-only estimator,
\[
\widetilde{Y}_{\mathrm{ot},i}
  \coloneqq \frac{A_i Y_i}{\widehat{\pi}_{ot}(\bX_i)}
  - \biggl(\frac{A_i}{\widehat{\pi}_{ot}(\bX_i)} - 1\biggr)
    \psi\bigl(\phi(\bX_i);\,\widehat{\balpha}_{ot}^{1,(-k(i))}\bigr),
\]
so that the centered IF satisfies $\widehat{\varphi}_{\mathrm{ot},i} = \widetilde{Y}_{\mathrm{ot},i} - \widehat{\mu}^1_{\mathrm{ot},k(i)}$.
Likewise, $\widehat{\zeta}_{j,i} = \psi(\phi(\bX_i);\widehat{\balpha}_{t,s_j}^{1,(-k(i))}) - \tildeE_{t,k(i)}[\psi(\cdot)]$.
All nuisance parameters $\widehat{\balpha}_{t,s_j}^{1,(-k)}$, $\widehat{\bgamma}_{s_j,1}^{(-k)}$, and $\widehat{\pi}_{ot}^{(-k)}$ are trained on data excluding the evaluation fold $\mathcal{D}_k$ (with the inner $M$-fold calibration/averaging described above), and $\widehat{\boldsymbol{\eta}}^{(-k)}$ is estimated via the inner $M$-fold loop within $\mathcal{T}_k = \mathcal{D}\setminus\mathcal{D}_k$.
Note that the source contribution is written in uncentered form; subtracting the fold-specific correction $\widehat{\delta}_{t,s_j}^{(-k)} = \tildeE_{s_j,k}[w(Y-\psi)]$ yields the centered influence component $\widehat{\xi}_{j,i}$ used in the variance-component derivations above.

The point estimate and variance are then derived from the \emph{same} set $\{\widehat{\Phi}_i\}_{i=1}^{N_{\mathrm{all}}}$:
\begin{align*}
\widehat{\mu}^1_{\mathrm{agg}} &= \bar{\Phi} \coloneqq \frac{1}{N_{\mathrm{all}}}\sum_{i=1}^{N_{\mathrm{all}}} \widehat{\Phi}_i, \\[4pt]
\widehat{V}_{\mathrm{agg}} &= \frac{1}{N_{\mathrm{all}}^2}\sum_{g\in\{t,s_1,\ldots,s_K\}}\;\sum_{i\in\mathcal{D}_g} \bigl(\widehat{\Phi}_i - \bar{\Phi}_g\bigr)^2,
\end{align*}
where $\bar{\Phi}_g \coloneqq \frac{1}{N_g}\sum_{i\in\mathcal{D}_g}\widehat{\Phi}_i$ is the \emph{within-group} mean for group $g\in\{t,s_1,\ldots,s_K\}$.
In the ANOVA decomposition $\text{TSS}=\text{WSS}+\text{BSS}$, the estimator $\widehat{V}_{\mathrm{agg}}$ uses only the within-group sum of squares $\text{WSS}/N_{\mathrm{all}}^2$; the between-group sum of squares $\text{BSS}/N_{\mathrm{all}}^2$ is excluded.

\textbf{Why within-group centering is necessary.}
Because $\widehat{\Phi}_i$ is scaled by $N_{\mathrm{all}}/N_g$ (i.e., $N_{\mathrm{all}}/N_t$ for target, $N_{\mathrm{all}}/N_{s_j}$ for source), the conditional expectations differ across groups:
$\E[\widehat{\Phi}_i\mid i\in\mathcal{D}_t]\approx (N_{\mathrm{all}}/N_t)\mu^1$ and
$\E[\widehat{\Phi}_i\mid i\in\mathcal{D}_{s_j}]\approx 0$ (when $\delta^*_{t,s_j}=0$).
The inferential target here is
\[
\frac{1}{N_{\mathrm{all}}^2}\sum_{i=1}^{N_{\mathrm{all}}}\Var(\Phi_i),
\]
rather than $\Var(\Phi)$ under a pooled-mixture distribution.
Global centering at $\bar{\Phi}$ estimates the latter object and introduces an additional between-group term
\[
\sum_g \rho_g\,(\mu_g-\mu)^2,
\]
equivalently the ANOVA component $\text{BSS}/N_{\mathrm{all}}^2$, which is of the same order as $V^*_{\mathrm{agg}}/N_{\mathrm{all}}$ and therefore does not vanish.
By contrast, $\text{WSS}/N_{\mathrm{all}}^2$ expands as
\[
\widehat{V}_{\mathrm{agg}}
= \frac{1}{N_t}\,\widehat{\sigma}^2_{T}
  + \sum_{j=1}^{K}\frac{1}{N_{s_j}}\,\widehat{\eta}_j^2\,\widehat{\sigma}^2_{S_j},
\]
where
\[
\widehat{\sigma}^2_T
=\widehat{\Var}_t\!\left((1-\sum_j\widehat{\eta}_j)\widehat{\varphi}_{\mathrm{ot},i}+\sum_j\widehat{\eta}_j\widehat{\zeta}_{j,i}\right).
\]
Hence $\widehat{\sigma}^2_T$ is not a target-only variance term; it already contains the shared-target cross-covariances (including terms corresponding to $\Cov(\widehat{\varphi}_{\mathrm{ot}},\widehat{\zeta}_j)$ and $\Cov(\widehat{\zeta}_j,\widehat{\zeta}_k)$ for $j\neq k$), fully consistent with the $\widehat{C}_{ot,s_j}$ and $\widehat{C}_{t,s_j,s_k}$ terms in Equation~\eqref{eq:final_opt}.

With balanced folds ($|\mathcal{D}_{t,k}|=N_t/K_f$, $|\mathcal{D}_{s_j,k}|=N_{s_j}/K_f$), $\bar{\Phi}$ coincides with the fold-average $(1/K_f)\sum_k \widehat{\mu}^1_{\mathrm{agg},k}$; with minor fold imbalance the difference is $O(N_{\mathrm{all}}^{-1})$.

This unified form ensures $\widehat{\mu}^1_{\mathrm{agg}} - \mu^1 = \frac{1}{N_{\mathrm{all}}}\sum_{i=1}^{N_{\mathrm{all}}}(\widehat{\Phi}_i - \mu^1)$, so the DML representation of the estimation error as an IF average is preserved. By the same arguments as above (nuisance consistency, $L_2$-Lipschitz continuity, and $\widehat{\boldsymbol{\eta}}^{(-k)}=\boldsymbol{\eta}^*+o_p(1)$), $\widehat{V}_{\mathrm{agg}} \xrightarrow{p} V^*_{\mathrm{agg}}$.
\end{proof}

\subsection{Technical Lemmas}

\begin{lemma}[CLT for IF Average]\label{lem:clt}
Under Assumptions \ref{asmp:data} and \ref{asmp:overlap}, the influence function average is asymptotically normal:
\[
\frac{1}{\sqrt{N}}\sum_{i=1}^{N}\varphi_{t,s}(D_i;\vartheta^*)
\;\xrightarrow{d}\; \mathcal{N}(0,V^*_{t,s}).
\]
\end{lemma}
\begin{proof}
The observations consist of two independent samples: $\{D_i: i \in \mathcal{D}_t\}$ drawn i.i.d.\ from the target distribution and $\{D_i: i \in \mathcal{D}_s\}$ drawn i.i.d.\ from the source distribution, with $N_t/N \to \rho_t \in (0,1)$ and $N_s/N \to \rho_s$.
Thus the summands $\{\varphi_{t,s}(D_i;\vartheta^*)\}_{i=1}^N$ are independent (but not identically distributed), have mean zero, and have finite $(2+\delta)$ moments by Assumption~\ref{asmp:overlap}(b). In particular, the Lindeberg condition holds.
We verify a Lyapunov condition, which implies the Lindeberg condition for this independent triangular array.
Let $X_{i,N}\coloneqq \varphi_{t,s}(D_i;\vartheta^*)$ so that $\E[X_{i,N}]=0$ and $\sum_{i=1}^{N}\Var(X_{i,N})=NV^*_{t,s}+o(N)$.
By Assumption \ref{asmp:overlap}(b), $\sup_{i,N}\E[|X_{i,N}|^{2+\delta}]<\infty$ for some $\delta>0$.
Therefore,
\[
\frac{1}{(NV^*_{t,s})^{1+\delta/2}}\sum_{i=1}^{N}\E[|X_{i,N}|^{2+\delta}]
\le
\frac{C N}{N^{1+\delta/2}}
=
CN^{-\delta/2}\to 0,
\]
which is the Lyapunov condition. Hence the Lindeberg--Feller CLT yields
\(
\frac{1}{\sqrt{N}}\sum_{i=1}^{N}X_{i,N}\xrightarrow{d}\mathcal{N}(0,V^*_{t,s})
\),
proving the lemma.
\end{proof}

\subsection{Extension to the Control Mean and the ATE}\label{sec:ate_extension}

The results above were stated for the treated potential outcome mean $\mu^1=\E_t[Y(1)]$ for notational simplicity. The same proof strategy applies to the control mean
\[
\mu^0 \coloneqq \E_t[Y(0)],
\]
by replacing the treatment level $1$ by $0$ throughout (including the corresponding nuisance parameters $(\boldsymbol{\alpha}^0,\boldsymbol{\gamma}_0)$ and weights).
In particular, define the control-arm weight
\[
w^{0}(\boldsymbol{X};\boldsymbol{\gamma}_0)\coloneqq \frac{I(A=0)}{e^{\phi(\boldsymbol{X})^T \boldsymbol{\gamma}_0}},
\]
and the cross-fitted estimator
\[
\widehat{\mu}^{0,\mathrm{CF}}_{t,s}
= \frac{1}{K_f}\sum_{k=1}^{K_f}
\Big\{
    \tildeE_{t,k}\big[\psi(\phi(\boldsymbol{X}); \widehat{\boldsymbol{\alpha}}_{t,s}^{0,(-k)})\big]
    + \tildeE_{s,k}\!\left[
        w^{0}(\boldsymbol{X};\widehat{\boldsymbol{\gamma}}_{s,0}^{(-k)})
        \big(Y - \psi(\phi(\boldsymbol{X}); \widehat{\boldsymbol{\alpha}}_{t,s}^{0,(-k)})\big)
    \right]
\Big\}.
\]
Let $\vartheta^{0,*}=(\boldsymbol{\alpha}^{0,*},\boldsymbol{\gamma}_0^*)$ denote the population targets of the control-arm calibrated losses (the same construction as for $\mu^1$ but with $I(A=0)$), and define the corresponding centered influence function $\varphi^{0}_{t,s}(D;\vartheta^{0,*})$ analogously.

\begin{theorem}[Asymptotic Normality of $\widehat{\mu}^{0,\mathrm{CF}}_{t,s}$]\label{thm:main_mu0}
Under the analogues of Assumptions \ref{asmp:data}--\ref{asmp:rsc} for the control arm, and assuming identification of $\mu^0(\vartheta^{0,*})=\mu^0$ under outcome-model or weight-model correctness (control-arm analogue of Lemma \ref{lem:ident_dr_mu}), if the total sparsity satisfies $\frac{s_{\mathrm{tot}}^2\log^2(p)}{N}\to 0$, then
\[
\sqrt{N}\,\big(\widehat{\mu}^{0,\mathrm{CF}}_{t,s} - \mu^0\big)
\;\xrightarrow{d}\;
\mathcal{N}\!\big(0,\,V^{0,*}_{t,s}\big),
\qquad
V^{0,*}_{t,s}\coloneqq \Var\!\big(\varphi^{0}_{t,s}(D;\vartheta^{0,*})\big).
\]
\end{theorem}

\begin{proof}
The proof is identical to that of Theorem \ref{thm:main} after replacing the treatment level $1$ by $0$ in the definition of the score contribution, the calibrated moment conditions, and the corresponding nuisance estimators.
\end{proof}

\begin{corollary}[Asymptotic Normality of the ATE]\label{cor:ate}
Let $\tau\coloneqq \mu^1-\mu^0$ and define the cross-fitted ATE estimator
\[
\widehat{\tau}^{\mathrm{CF}}_{t,s}\coloneqq \widehat{\mu}^{1,\mathrm{CF}}_{t,s}-\widehat{\mu}^{0,\mathrm{CF}}_{t,s}.
\]
Under the conditions of Theorems \ref{thm:main} and \ref{thm:main_mu0}, we have the asymptotic linear representation
\[
\sqrt{N}\,\big(\widehat{\tau}^{\mathrm{CF}}_{t,s}-\tau\big)
=
\frac{1}{\sqrt{N}}\sum_{i=1}^N
\Big(
\varphi_{t,s}(D_i;\vartheta^*)
-
\varphi^{0}_{t,s}(D_i;\vartheta^{0,*})
\Big)
+o_p(1),
\]
and hence
\[
\sqrt{N}\,\big(\widehat{\tau}^{\mathrm{CF}}_{t,s}-\tau\big)
\;\xrightarrow{d}\;
\mathcal{N}\!\big(0,\,V^{\tau,*}_{t,s}\big),
\qquad
V^{\tau,*}_{t,s}\coloneqq \Var\!\Big(\varphi_{t,s}(D;\vartheta^*)-\varphi^{0}_{t,s}(D;\vartheta^{0,*})\Big).
\]
\end{corollary}

\begin{lemma}[Cross-fitted $\ell_1$-penalized nuisance rates (adapted from \citet{hou2025efficient}, Lemma A21)]\label{lem:nuisance_rates}
Fix a constant number of folds $K_f$. Let the full sample $\mathcal{D}=\{D_1,\ldots,D_N\}$ be partitioned into folds $\mathcal{D}=\bigcup_{k=1}^{K_f}\mathcal{D}_k$, and let $\mathcal{D}^{(-k)}=\mathcal{D}\setminus\mathcal{D}_k$ denote the out-of-fold data. For each fold $k$, let $\widehat{\boldsymbol{\gamma}}^{(-k)}$ be an estimator trained on $\mathcal{D}^{(-k)}$ (hence independent of $\mathcal{D}_k$).
For an observation index $i\in\{1,\ldots,N\}$, let $\mathrm{fold}(i)\in\{1,\ldots,K_f\}$ denote the fold index such that $D_i\in\mathcal{D}_{\mathrm{fold}(i)}$; we use $\widehat{\boldsymbol{\gamma}}^{(-\mathrm{fold}(i))}$ to denote the out-of-fold nuisance estimate used when evaluating the loss contribution from $D_i$.
Consider the cross-fitted $\ell_1$-penalized estimator
\[
\widehat{\boldsymbol{\beta}}
\;\in\;
\arg\min_{\boldsymbol{\beta}\in\R^p}
\left\{
\frac{1}{K_f}\sum_{k=1}^{K_f}\ell_k(\boldsymbol{\beta};\widehat{\boldsymbol{\gamma}}^{(-k)})
\;+\;
\lambda\|\boldsymbol{\beta}\|_1
\right\},
\qquad
\ell_k(\boldsymbol{\beta};\boldsymbol{\gamma})
\coloneqq
\tildeE_k\!\big[\ell(D;\boldsymbol{\beta},\boldsymbol{\gamma})\big],
\]
where $\ell(D;\boldsymbol{\beta},\boldsymbol{\gamma})$ is convex in $\boldsymbol{\beta}$.
Define also the fold-averaged empirical risk
\[
L_N(\boldsymbol{\beta}) \coloneqq \frac{1}{K_f}\sum_{k=1}^{K_f}\ell_k(\boldsymbol{\beta};\widehat{\boldsymbol{\gamma}}^{(-k)}).
\]
Let $\overline{\boldsymbol{\beta}}$ be the population minimizer at a population nuisance target $\overline{\boldsymbol{\gamma}}$,
\[
\overline{\boldsymbol{\beta}}\in\arg\min_{\boldsymbol{\beta}\in\R^p}\E[\ell(D;\boldsymbol{\beta},\overline{\boldsymbol{\gamma}})],
\qquad s_\beta\coloneqq\|\overline{\boldsymbol{\beta}}\|_0.
\]
Assume the following high-dimensional regularity conditions (cf.\ Assumption 6 of \citet{hou2025efficient} and the RSC framework of \citet{negahban2012unified}):
\begin{enumerate}
    \item[(a)] (Sub-Gaussian covariates) Each coordinate of $\phi(\boldsymbol{X})$ is sub-Gaussian with $\max_{1\le j\le p}\|\phi_j(\boldsymbol{X})\|_{\psi_2}\le K_\phi$ (allowing $\|\phi(\boldsymbol{X})\|_\infty$ to grow at the usual $\sqrt{\log p}$ rate).
    \item[(b)] (Response/link tails and smoothness) $Y$ is uniformly sub-Gaussian (or bounded), and the GLM mean function and its derivative are uniformly bounded on a neighborhood of $\overline{\boldsymbol{\beta}}$.
    \item[(c)] (Identifiability) The population Gram matrix is non-degenerate: $\lambda_{\min}(\E[\phi(\boldsymbol{X})\phi(\boldsymbol{X})^T])\ge c>0$.
    \item[(d)] (Weight regularity) Any weights induced by $\boldsymbol{\gamma}$ (e.g., density-ratio type weights) have finite moments and are locally Lipschitz in $\boldsymbol{\gamma}$ on a neighborhood of $\overline{\boldsymbol{\gamma}}$ (sufficient for the score stability bound used below).
    \item[(e)] (Restricted strong convexity) With probability tending to 1, $L_N(\cdot)$ satisfies an RSC inequality on the standard cone associated with $\mathrm{supp}(\overline{\boldsymbol{\beta}})$.
\end{enumerate}
In our FACE-HD setting, we impose these regularity conditions in addition to Assumptions \ref{asmp:data}, \ref{asmp:overlap}, and \ref{asmp:rsc} (noting that (a) matches Assumption \ref{asmp:data}(a), while (d) is ensured by the overlap/weight-moment assumptions and local smoothness). Choose $\lambda\asymp\sqrt{\log(p)/N}$. Then:
\begin{remark}[On the choice of $N$ in $\lambda$]
The scaling $\lambda\asymp\sqrt{\log(p)/N}$ is stated in terms of the total sample size $N$ for notational simplicity. In applications where the nuisance loss for a given arm/site pair is effectively averaged over an arm-specific subsample (e.g., treated units in the source), one may equivalently write $\lambda\asymp\sqrt{\log(p)/N_{\mathrm{eff}}}$ where $N_{\mathrm{eff}}$ is the corresponding effective sample size. Under the sampling-fraction regime and overlap/positivity conditions stated above, $N_{\mathrm{eff}}\asymp N$, so the rate statements are unchanged.
\end{remark}
\begin{enumerate}
    \item \textbf{(General)} 
    \[
    \|\widehat{\boldsymbol{\beta}}-\overline{\boldsymbol{\beta}}\|_2
    =
    O_p\!\left(\sqrt{\frac{s_\beta\log(p)}{N}}\right)
    \;+\;
    O_p\!\left(\sup_{1\le k\le K_f}\|\widehat{\boldsymbol{\gamma}}^{(-k)}-\overline{\boldsymbol{\gamma}}\|_2\right).
    \]
    \item \textbf{(Special)} If, for each fold $k$, the conditional foldwise score at $\overline{\boldsymbol{\beta}}$ satisfies
    \[
    \E\!\left[\nabla_{\boldsymbol{\beta}}\ell(D;\overline{\boldsymbol{\beta}},\widehat{\boldsymbol{\gamma}}^{(-k)})\,\middle|\,\mathcal{D}^{(-k)}\right]
    =\mathbf{0},
    \]
    then the plug-in score term disappears and
    \[
    \|\widehat{\boldsymbol{\beta}}-\overline{\boldsymbol{\beta}}\|_2
    =
    O_p\!\left(\sqrt{\frac{s_\beta\log(p)}{N}}\right).
    \]
\end{enumerate}
Moreover, we can convert the $\ell_2$-rate to an $\ell_1$-rate by an explicit cone argument.
Let $\Delta\coloneqq \widehat{\boldsymbol{\beta}}-\overline{\boldsymbol{\beta}}$ and $\mathcal{J}\coloneqq \mathrm{supp}(\overline{\boldsymbol{\beta}})$ with $|\mathcal{J}|=s_\beta$.
From the proof below (see the KKT-based cone inequality leading to \eqref{eq:cone_ineq}), on a high-probability event we have
\[
\|\Delta_{\mathcal{J}^c}\|_1 \le 3\|\Delta_{\mathcal{J}}\|_1 + \frac{2}{\lambda}\|B\|_\infty\|\Delta\|_1,
\]
where $\|B\|_\infty$ is the plug-in score term induced by cross-fitting.
With $\lambda\asymp\sqrt{\log(p)/N}$ and $\|B\|_\infty=O_p(\sqrt{\log(p)/N})$ (proved below), we have $(2/\lambda)\|B\|_\infty\le 1/2$ w.p.a.1, and hence
\[
\|\Delta\|_1
=
\|\Delta_{\mathcal{J}}\|_1+\|\Delta_{\mathcal{J}^c}\|_1
\le
4\|\Delta_{\mathcal{J}}\|_1
\le
4\sqrt{s_\beta}\|\Delta\|_2.
\]
Combining this with the $\ell_2$ bounds above yields the corresponding $\ell_1$-rate:
\[
\|\widehat{\boldsymbol{\beta}}-\overline{\boldsymbol{\beta}}\|_1
=
O_p\!\left(s_\beta\sqrt{\frac{\log(p)}{N}}\right)
\;+\;
O_p\!\left(\sqrt{s_\beta}\sup_{1\le k\le K_f}\|\widehat{\boldsymbol{\gamma}}^{(-k)}-\overline{\boldsymbol{\gamma}}\|_2\right),
\]
with the second term dropping in the special case.

\textbf{Instantiation for FACE-HD calibrated losses.} In our application, $\boldsymbol{\beta}$ is instantiated as either $\boldsymbol{\alpha}^1$ or $\boldsymbol{\gamma}_1$. If Assumption \ref{asmp:rsc}(d) holds (as in the foldwise score-centering condition used by \citet{hou2025efficient}), then we may invoke the \textbf{(Special)} case above and obtain
\[
\lVert \widehat{\boldsymbol{\alpha}}^{1,(-k)}_{t,s} - \boldsymbol{\alpha}^{1,*}\rVert_q
= O_p\!\left(s_\alpha^{1-1/q}\sqrt{\frac{\log(p)}{N}}\right),
\qquad
\lVert \widehat{\boldsymbol{\gamma}}^{(-k)}_{s,1} - \boldsymbol{\gamma}^{*}_1\rVert_q
= O_p\!\left(s_\gamma^{1-1/q}\sqrt{\frac{\log(p)}{N}}\right),
\qquad q\in\{1,2\}.
\]
\end{lemma}
\begin{proof}
We give a detailed argument in our notation; it follows the structure of \citet{hou2025efficient} (proof of Lemma A21).
Write the (random) fold-averaged empirical loss as
\[
L_N(\boldsymbol{\beta}) \coloneqq \frac{1}{K_f}\sum_{k=1}^{K_f}\ell_k(\boldsymbol{\beta};\widehat{\boldsymbol{\gamma}}^{(-k)}),
\qquad
\widehat{\boldsymbol{\beta}}\in\arg\min_{\boldsymbol{\beta}}\{L_N(\boldsymbol{\beta})+\lambda\|\boldsymbol{\beta}\|_1\}.
\]
Let $\Delta\coloneqq \widehat{\boldsymbol{\beta}}-\overline{\boldsymbol{\beta}}$, and let $\mathcal{J}\coloneqq \mathrm{supp}(\overline{\boldsymbol{\beta}})$ with $|\mathcal{J}|=s_\beta$.

\textbf{Step 1 (Basic inequality).}
By optimality,
\[
L_N(\overline{\boldsymbol{\beta}}+\Delta)+\lambda\|\overline{\boldsymbol{\beta}}+\Delta\|_1
\le
L_N(\overline{\boldsymbol{\beta}})+\lambda\|\overline{\boldsymbol{\beta}}\|_1,
\]
so
\[
L_N(\overline{\boldsymbol{\beta}}+\Delta)-L_N(\overline{\boldsymbol{\beta}})
\le
\lambda(\|\overline{\boldsymbol{\beta}}\|_1-\|\overline{\boldsymbol{\beta}}+\Delta\|_1)
\le
\lambda(\|\Delta_{\mathcal{J}}\|_1-\|\Delta_{\mathcal{J}^c}\|_1).
\]

\textbf{Step 2 (RSC expansion).}
By convexity,
\[
L_N(\overline{\boldsymbol{\beta}}+\Delta)-L_N(\overline{\boldsymbol{\beta}})
\ge
\nabla L_N(\overline{\boldsymbol{\beta}})^T\Delta
\;+\;
\mathcal{R}_N(\Delta),
\]
where the remainder $\mathcal{R}_N(\Delta)$ satisfies the RSC lower bound (Assumption \ref{asmp:rsc}(a) as applied to this loss)
\[
\mathcal{R}_N(\Delta)
\ge
\kappa_1\|\Delta\|_2^2-\kappa_2\frac{\log(p)}{N}\|\Delta\|_1^2
\]
on the standard restricted cone. Combining with Step 1 yields
\begin{equation}\label{eq:nuis_prelim_bound}
\kappa_1\|\Delta\|_2^2
\le
\|\nabla L_N(\overline{\boldsymbol{\beta}})\|_\infty\|\Delta\|_1
\;+\;
\lambda(\|\Delta_{\mathcal{J}}\|_1-\|\Delta_{\mathcal{J}^c}\|_1)
\;+\;
\kappa_2\frac{\log(p)}{N}\|\Delta\|_1^2.
\end{equation}

\textbf{Step 3 (Score decomposition and choice of $\lambda$).}
Decompose the score
\[
\nabla L_N(\overline{\boldsymbol{\beta}})
=
\underbrace{\big(\nabla L_N(\overline{\boldsymbol{\beta}})-\E[\nabla L_N(\overline{\boldsymbol{\beta}})\mid \{\mathcal{D}^{(-k)}\}_{k=1}^{K_f}]\big)}_{\text{stochastic term }A}
\;+\;
\underbrace{\E[\nabla L_N(\overline{\boldsymbol{\beta}})\mid \{\mathcal{D}^{(-k)}\}_{k=1}^{K_f}]}_{\text{plug-in term }B}.
\]
By cross-fitting, conditional on $\{\mathcal{D}^{(-k)}\}$, each fold average $\tildeE_k[\nabla_{\boldsymbol{\beta}}\ell(D;\overline{\boldsymbol{\beta}},\widehat{\boldsymbol{\gamma}}^{(-k)})]$ is an average of independent sub-exponential random vectors with uniformly controlled coordinates, so a union bound with Bernstein/Hoeffding inequalities gives
\[
\|A\|_\infty = O_p\!\left(\sqrt{\frac{\log(p)}{N}}\right).
\]
Choose $\lambda\asymp\sqrt{\log(p)/N}$ so that $\lambda\ge 2\|A\|_\infty$ w.p.a.1.

For the plug-in term, note that
\[
B-\E[\nabla_{\boldsymbol{\beta}}\ell(D;\overline{\boldsymbol{\beta}},\overline{\boldsymbol{\gamma}})]
=
\E[\nabla_{\boldsymbol{\beta}}\ell(D;\overline{\boldsymbol{\beta}},\widehat{\boldsymbol{\gamma}}^{(-\mathrm{fold}(D))})
-\nabla_{\boldsymbol{\beta}}\ell(D;\overline{\boldsymbol{\beta}},\overline{\boldsymbol{\gamma}})
\mid \{\mathcal{D}^{(-k)}\}],
\]
and the weight-regularity/Lipschitz condition implies
\[
\|B-\E[\nabla_{\boldsymbol{\beta}}\ell(D;\overline{\boldsymbol{\beta}},\overline{\boldsymbol{\gamma}})]\|_2
\le
C\sup_{1\le k\le K_f}\|\widehat{\boldsymbol{\gamma}}^{(-k)}-\overline{\boldsymbol{\gamma}}\|_2.
\]
In the special case, $\E[\nabla_{\boldsymbol{\beta}}\ell(D;\overline{\boldsymbol{\beta}},\widehat{\boldsymbol{\gamma}}^{(-k)})\mid\mathcal{D}^{(-k)}]=0$ for all $k$, hence $B=\mathbf{0}$.

\textbf{Step 4 (Cone property and $\ell_2$ rate).}
Plugging $\|\nabla L_N(\overline{\boldsymbol{\beta}})\|_\infty\le \|A\|_\infty+\|B\|_\infty$ into \eqref{eq:nuis_prelim_bound} and using $\lambda\ge 2\|A\|_\infty$ yields the cone inequality
\begin{equation}\label{eq:cone_ineq}
\|\Delta_{\mathcal{J}^c}\|_1 \le 3\|\Delta_{\mathcal{J}}\|_1 + \frac{2}{\lambda}\|B\|_\infty\|\Delta\|_1.
\end{equation}
On the event $\{(2/\lambda)\|B\|_\infty\le 1/2\}$, \eqref{eq:cone_ineq} implies $\|\Delta_{\mathcal{J}^c}\|_1\le 7\|\Delta_{\mathcal{J}}\|_1$ and hence $\|\Delta\|_1\le 8\|\Delta_{\mathcal{J}}\|_1\le 8\sqrt{s_\beta}\|\Delta\|_2$.
Plugging this back into the RSC lower bound (Assumption \ref{asmp:rsc}(a) as applied to this loss) absorbs the $\kappa_2(\log(p)/N)\|\Delta\|_1^2$ term and yields
\[
\|\Delta\|_2
=
O_p(\sqrt{s_\beta}\lambda)
\;+\;
O_p\!\left(\sup_{1\le k\le K_f}\|\widehat{\boldsymbol{\gamma}}^{(-k)}-\overline{\boldsymbol{\gamma}}\|_2\right).
\]
This is the general-case $\ell_2$ rate. In the special case, $B=\mathbf{0}$ and the second term drops.

\textbf{Step 5 ($\ell_1$ rate).}
From the cone property and $\|\Delta_{\mathcal{J}}\|_1\le \sqrt{s_\beta}\|\Delta\|_2$,
\[
\|\Delta\|_1
=
O_p(s_\beta\lambda)
\;+\;
O_p\!\left(\sqrt{s_\beta}\sup_{1\le k\le K_f}\|\widehat{\boldsymbol{\gamma}}^{(-k)}-\overline{\boldsymbol{\gamma}}\|_2\right),
\]
with the second term dropping in the special case.

Finally, for FACE-HD, the \textbf{(Special)} conclusion may be invoked for $\boldsymbol{\beta}=\boldsymbol{\alpha}^1$ and $\boldsymbol{\beta}=\boldsymbol{\gamma}_1$ only when the foldwise score-centering condition is assumed/verified as in Assumption \ref{asmp:rsc}(d). Otherwise, one should use the \textbf{(General)} bounds, which explicitly propagate the plug-in nuisance error.
\end{proof}

\begin{lemma}[Orthogonality from Calibrated Losses]\label{lem:ortho_foc}
Under Assumption \ref{asmp:rsc}(c), let $\vartheta^* = (\boldsymbol{\alpha}^{1,*}, \boldsymbol{\gamma}_1^*)$ be the population targets of the calibrated losses in the two-level algorithm (main manuscript appendix, Section~\ref{app:two_level_cf}; Equations~\eqref{eq:gamma_calibrated_loss} and \eqref{eq:alpha_calibrated_loss}):
\begin{align*}
L_{\gamma}(\boldsymbol{\gamma}; \boldsymbol{\alpha}^{1,*}) &= \E_t \left[ \nabla_{\boldsymbol{\alpha}} \psi(\phi(\boldsymbol{X}); \boldsymbol{\alpha}^{1,*})^T \boldsymbol{\gamma} \right] + \E_{s} \left[ I(A=1) e^{-\phi(\boldsymbol{X})^T \boldsymbol{\gamma}} \psi'(\mathcal{T}(\phi(\boldsymbol{X})^T \boldsymbol{\alpha}^{1,*})) \right] \\
L_{\alpha}(\boldsymbol{\alpha}; \boldsymbol{\gamma}_1^*) &= \E_{s} \left[ w_\tau(\boldsymbol{X}; \boldsymbol{\gamma}_1^*) \left( -Y \cdot (\phi(\boldsymbol{X})^T \boldsymbol{\alpha}) + \Psi(\phi(\boldsymbol{X})^T \boldsymbol{\alpha}) \right) \right]
\end{align*}
(Ignoring regularization terms, which vanish in population). For canonical-link GLMs,
\[
\nabla_{\boldsymbol{\alpha}} \psi(\phi(\boldsymbol{X}); \boldsymbol{\alpha})
= \phi(\boldsymbol{X})\,\psi'(\phi(\boldsymbol{X})^T\boldsymbol{\alpha}).
\]
\begin{remark}[Penalization and approximate orthogonality]
The population moment conditions below are derived from the unpenalized population losses. In finite samples, the $\ell_1$ penalty induces a KKT subgradient term of order $\lambda$ in the estimating equations. Under the nuisance rates in Lemma~\ref{lem:nuisance_rates} and the sparsity scaling used in Theorem~\ref{thm:main}, this penalty-induced perturbation contributes at most an $o_p(N^{-1/2})$ term to the final DR remainder and therefore does not affect the first-order limit theory.
\end{remark}
Under asymptotically inactive truncation (Assumption~\ref{asmp:overlap}(c)), the population first-order optimality conditions are therefore equivalent to the following moment conditions at $\vartheta^*$:
\begin{align*}
(i) \quad & \E_t\big[\nabla_{\boldsymbol{\alpha}} \psi(\phi(\boldsymbol{X}); \boldsymbol{\alpha}^{1,*})\big]
- \E_s\left[
    w(\boldsymbol{X};\boldsymbol{\gamma}_1^*)
    \nabla_{\boldsymbol{\alpha}} \psi(\phi(\boldsymbol{X}); \boldsymbol{\alpha}^{1,*})
\right] = \mathbf{0} \\
(ii) \quad & \E_s\left[
    w(\boldsymbol{X};\boldsymbol{\gamma}_1^*)
    \phi(\boldsymbol{X})
    \{Y - \psi(\phi(\boldsymbol{X}); \boldsymbol{\alpha}^{1,*})\}
\right] = \mathbf{0}.
\end{align*}
\end{lemma}

\begin{proof}
We derive the moment conditions by taking the gradients of the population loss functions $L_\gamma$ and $L_\alpha$ with respect to the parameter being optimized and setting them to zero at the minimum $\vartheta^*$.

\textbf{Condition (i) from FOC of $L_{\gamma}(\boldsymbol{\gamma}; \boldsymbol{\alpha}^{1,*})$ w.r.t. $\boldsymbol{\gamma}$:}
The gradient of $L_{\gamma}(\boldsymbol{\gamma}; \boldsymbol{\alpha}^{1,*})$ with respect to $\boldsymbol{\gamma}$ is:
\[
\begin{aligned}
\nabla_{\boldsymbol{\gamma}} L_{\gamma}(\boldsymbol{\gamma}; \boldsymbol{\alpha}^{1,*})
&= \E_t \left[ \nabla_{\boldsymbol{\alpha}} \psi(\phi(\boldsymbol{X}); \boldsymbol{\alpha}^{1,*}) \right] \\
&\quad + \E_{s} \left[ I(A=1) (-\phi(\boldsymbol{X})) e^{-\phi(\boldsymbol{X})^T \boldsymbol{\gamma}} \psi'(\mathcal{T}(\phi(\boldsymbol{X})^T \boldsymbol{\alpha}^{1,*})) \right]
\end{aligned}
\]
Setting this gradient to zero at $\boldsymbol{\gamma} = \boldsymbol{\gamma}_1^*$ gives:
\[
\begin{aligned}
\E_t \left[ \nabla_{\boldsymbol{\alpha}} \psi(\phi(\boldsymbol{X}); \boldsymbol{\alpha}^{1,*}) \right]
&= \E_{s} \left[ \phi(\boldsymbol{X}) w(\boldsymbol{X}; \boldsymbol{\gamma}_1^*) \psi'(\mathcal{T}(\phi(\boldsymbol{X})^T \boldsymbol{\alpha}^{1,*})) \right].
\end{aligned}
\]
This equation is constructed to enforce the moment condition (i). Note that for standard GLMs with canonical links,
\[
\nabla_{\boldsymbol{\alpha}} \psi(\phi(\boldsymbol{X}); \boldsymbol{\alpha}) = \phi(\boldsymbol{X}) \psi'(\phi(\boldsymbol{X})^T \boldsymbol{\alpha}).
\]

\textbf{Truncation Error Analysis:} Define the event $\mathcal{E}_{\alpha,N} = \{|\phi(\boldsymbol{X})^T \boldsymbol{\alpha}^{1,*}| \le M_{\tau,N}\}$. By Assumption~\ref{asmp:overlap}(c), $P(\mathcal{E}_{\alpha,N}^c \mid R=s)=o(N^{-1})$. On $\mathcal{E}_{\alpha,N}$, $\mathcal{T}(\phi(\boldsymbol{X})^T \boldsymbol{\alpha}^{1,*}) = \phi(\boldsymbol{X})^T \boldsymbol{\alpha}^{1,*}$, so $\psi'(\mathcal{T}(\cdot)) = \psi'(\cdot)$ exactly. Therefore, replacing $\psi'(\mathcal{T}(\phi(\boldsymbol{X})^T \boldsymbol{\alpha}^{1,*}))$ by $\psi'(\phi(\boldsymbol{X})^T \boldsymbol{\alpha}^{1,*})$ in the population FOC changes the source-side moment by at most an $o(1)$ term, and condition (i) holds asymptotically.

\textbf{Condition (ii) from FOC of $L_{\alpha}(\boldsymbol{\alpha}; \boldsymbol{\gamma}_1^*)$ w.r.t. $\boldsymbol{\alpha}$:}
The loss $L_{\alpha}$ is the population weighted GLM log-likelihood. Its gradient with respect to $\boldsymbol{\alpha}$ is:
\[
\nabla_{\boldsymbol{\alpha}} L_{\alpha}(\boldsymbol{\alpha}; \boldsymbol{\gamma}_1^*) = \E_{s} \left[ w_\tau(\boldsymbol{X}; \boldsymbol{\gamma}_1^*) \left( -Y \cdot \phi(\boldsymbol{X}) + \nabla_{\boldsymbol{\alpha}}\Psi(\phi(\boldsymbol{X})^T \boldsymbol{\alpha}) \right) \right]
\]
Using the GLM identity $\nabla_{\boldsymbol{\alpha}}\Psi(u) = \psi(u) \phi(\boldsymbol{X})$ where $u = \phi(\boldsymbol{X})^T \boldsymbol{\alpha}$:
\[
\nabla_{\boldsymbol{\alpha}} L_{\alpha}(\boldsymbol{\alpha}; \boldsymbol{\gamma}_1^*) = \E_{s} \left[ w_\tau(\boldsymbol{X}; \boldsymbol{\gamma}_1^*) \phi(\boldsymbol{X}) \left( \psi(\phi(\boldsymbol{X})^T \boldsymbol{\alpha}) - Y \right) \right]
\]
Setting this gradient to zero at $\boldsymbol{\alpha} = \boldsymbol{\alpha}^{1,*}$ gives:
\[
\E_{s} \left[ w_\tau(\boldsymbol{X}; \boldsymbol{\gamma}_1^*) \phi(\boldsymbol{X}) \left( Y - \psi(\phi(\boldsymbol{X})^T \boldsymbol{\alpha}^{1,*}) \right) \right] = \mathbf{0}.
\]
Finally, under Assumption~\ref{asmp:overlap}(c), truncation is asymptotically inactive at $\boldsymbol{\gamma}_1^*$, so $w_\tau(\boldsymbol{X};\boldsymbol{\gamma}_1^*) = w(\boldsymbol{X};\boldsymbol{\gamma}_1^*)$ w.p.a.1. This yields condition (ii).
\end{proof}

\begin{lemma}[Bias-Variance Bounds]\label{lem:bias_bound}
Let $\vartheta_n = (\boldsymbol{\alpha}_n^1, \boldsymbol{\gamma}_{1,n})$ be a sequence of nuisance parameters converging to $\vartheta^*=(\boldsymbol{\alpha}^{1,*},\boldsymbol{\gamma}_1^*)$.
\begin{enumerate}
    \item \textbf{(Bias Bound - Double Robustness)} If $\vartheta^*$ minimizes the population calibrated losses (Lemma \ref{lem:ortho_foc} holds), then:
    \[
    \begin{aligned}
    \lvert \E[\widetilde{\varphi}_{t,s}(D; \vartheta_n)] - \E[\widetilde{\varphi}_{t,s}(D; \vartheta^*)] \rvert
    &\le C \cdot \left(
        \lVert \boldsymbol{\alpha}_n^1 - \boldsymbol{\alpha}^{1,*} \rVert_2^2
        +
        \lVert \boldsymbol{\gamma}_{1,n} - \boldsymbol{\gamma}_1^* \rVert_2^2 \right. \\
    &\quad\left. +
        \lVert \boldsymbol{\alpha}_n^1 - \boldsymbol{\alpha}^{1,*} \rVert_2 \cdot \lVert \boldsymbol{\gamma}_{1,n} - \boldsymbol{\gamma}_1^* \rVert_2
    \right)
    \end{aligned}
    \]
    \item \textbf{(Variance Bound)}
    \[
    \Var(\widetilde{\varphi}_{t,s}(D; \vartheta_n) - \widetilde{\varphi}_{t,s}(D; \vartheta^*))
    \le C \cdot \left( \lVert \boldsymbol{\alpha}_n^1 - \boldsymbol{\alpha}^{1,*} \rVert_2^2 + \lVert \boldsymbol{\gamma}_{1,n} - \boldsymbol{\gamma}_1^* \rVert_2^2 \right)
    \]
\end{enumerate}
\end{lemma}
\begin{proof}
We analyze the bias term $\E[\widetilde{\varphi}_{t,s}(D; \vartheta_n) - \widetilde{\varphi}_{t,s}(D; \vartheta^*)]$.
Let $\Delta_\alpha = \boldsymbol{\alpha}_n^1 - \boldsymbol{\alpha}^{1,*}$ and $\Delta_\gamma = \boldsymbol{\gamma}_{1,n} - \boldsymbol{\gamma}_1^*$.
We decompose the bias using an intermediate step:
\[
\text{Bias} = \underbrace{\E[\widetilde{\varphi}_{t,s}(D; \boldsymbol{\alpha}_n^1, \boldsymbol{\gamma}_{1,n}) - \widetilde{\varphi}_{t,s}(D; \boldsymbol{\alpha}^{1,*}, \boldsymbol{\gamma}_{1,n})]}_{\text{Bias from } \boldsymbol{\alpha}}
+ \underbrace{\E[\widetilde{\varphi}_{t,s}(D; \boldsymbol{\alpha}^{1,*}, \boldsymbol{\gamma}_{1,n}) - \widetilde{\varphi}_{t,s}(D; \boldsymbol{\alpha}^{1,*}, \boldsymbol{\gamma}_1^*)]}_{\text{Bias from } \boldsymbol{\gamma}}
\]

\textbf{Step 1 (Bias from $\boldsymbol{\alpha}$: mean-value expansion + orthogonality).}
Define the path $h_\alpha(t)\coloneqq \E[\widetilde{\varphi}_{t,s}(D;\boldsymbol{\alpha}^{1,*}+t\Delta_\alpha,\boldsymbol{\gamma}_{1,n})]$ for $t\in[0,1]$.
By the fundamental theorem of calculus (and dominated convergence, justified by Assumption \ref{asmp:overlap}(b) and local boundedness of the derivatives), we have
\[
h_\alpha(1)-h_\alpha(0)
=
\int_0^1 h_\alpha'(t)\,dt
=
\int_0^1 \E\!\left[\nabla_{\boldsymbol{\alpha}}\widetilde{\varphi}_{t,s}(D;\boldsymbol{\alpha}^{1,*}+t\Delta_\alpha,\boldsymbol{\gamma}_{1,n})\right]^T\Delta_\alpha\,dt.
\]
Hence
\[
\text{Bias from } \boldsymbol{\alpha}
=
\left(\E\!\left[\nabla_{\boldsymbol{\alpha}}\widetilde{\varphi}_{t,s}(D;\boldsymbol{\alpha}^{1,*},\boldsymbol{\gamma}_1^*)\right]\right)^T\Delta_\alpha
\;+\;
\int_0^1 \left\{\E\!\left[\nabla_{\boldsymbol{\alpha}}\widetilde{\varphi}_{t,s}(D;\boldsymbol{\alpha}^{1,*}+t\Delta_\alpha,\boldsymbol{\gamma}_{1,n})\right]-\E\!\left[\nabla_{\boldsymbol{\alpha}}\widetilde{\varphi}_{t,s}(D;\boldsymbol{\alpha}^{1,*},\boldsymbol{\gamma}_1^*)\right]\right\}^T\Delta_\alpha\,dt.
\]
The first term is zero by Lemma \ref{lem:ortho_foc}(i). For the second term, we use Cauchy--Schwarz:
\[
\left|\text{Bias from } \boldsymbol{\alpha}\right|
\le
\|\Delta_\alpha\|_2 \int_0^1
\left\|
\E\!\left[\nabla_{\boldsymbol{\alpha}}\widetilde{\varphi}_{t,s}(D;\boldsymbol{\alpha}^{1,*}+t\Delta_\alpha,\boldsymbol{\gamma}_{1,n})\right]-\E\!\left[\nabla_{\boldsymbol{\alpha}}\widetilde{\varphi}_{t,s}(D;\boldsymbol{\alpha}^{1,*},\boldsymbol{\gamma}_1^*)\right]
\right\|_2\,dt.
\]
The mapping $\vartheta\mapsto \E[\nabla_{\boldsymbol{\alpha}}\widetilde{\varphi}_{t,s}(D;\vartheta)]$ is locally Lipschitz in $\vartheta$ by the explicit derivative formula below, the bounded-slope/Lipschitz condition on $\psi'$ (Assumption \ref{asmp:overlap}(b)(i)), and the bounded-weights/moment conditions (Assumption \ref{asmp:overlap}(a,b)). Therefore, for $\vartheta_n$ sufficiently close to $\vartheta^*$,
\[
\left\|
\E\!\left[\nabla_{\boldsymbol{\alpha}}\widetilde{\varphi}_{t,s}(D;\boldsymbol{\alpha}^{1,*}+t\Delta_\alpha,\boldsymbol{\gamma}_{1,n})\right]-\E\!\left[\nabla_{\boldsymbol{\alpha}}\widetilde{\varphi}_{t,s}(D;\boldsymbol{\alpha}^{1,*},\boldsymbol{\gamma}_1^*)\right]
\right\|_2
\le
C\left(\|\Delta_\alpha\|_2+\|\Delta_\gamma\|_2\right),
\]
uniformly over $t\in[0,1]$, and hence
\[
\left|\text{Bias from } \boldsymbol{\alpha}\right|
\le
C\left(\|\Delta_\alpha\|_2^2+\|\Delta_\alpha\|_2\|\Delta_\gamma\|_2\right).
\]

\textbf{Step 2 (Bias from $\boldsymbol{\gamma}$: mean-value expansion + orthogonality).}
Similarly define $h_\gamma(t)\coloneqq \E[\widetilde{\varphi}_{t,s}(D;\boldsymbol{\alpha}^{1,*},\boldsymbol{\gamma}_1^*+t\Delta_\gamma)]$.
Then
\[
\text{Bias from } \boldsymbol{\gamma}
=
\int_0^1 \E\!\left[\nabla_{\boldsymbol{\gamma}}\widetilde{\varphi}_{t,s}(D;\boldsymbol{\alpha}^{1,*},\boldsymbol{\gamma}_1^*+t\Delta_\gamma)\right]^T\Delta_\gamma\,dt.
\]
Adding and subtracting $\E[\nabla_{\boldsymbol{\gamma}}\widetilde{\varphi}_{t,s}(D;\boldsymbol{\alpha}^{1,*},\boldsymbol{\gamma}_1^*)]$, the leading term is zero by Lemma \ref{lem:ortho_foc}(ii), and the remainder is bounded (again by local Lipschitz continuity and Cauchy--Schwarz) as
\[
\left|\text{Bias from } \boldsymbol{\gamma}\right|
\le
C\|\Delta_\gamma\|_2^2.
\]

\textbf{Step 3 (Collecting).}
Combining the two displays yields the stated second-order bias bound.

\textbf{Derivative identities used above.}
For example, the gradient w.r.t. $\boldsymbol{\alpha}$ at $(\boldsymbol{\alpha}^{1,*},\boldsymbol{\gamma}_1^*)$ is:
\[
\begin{aligned}
\E\left[ \nabla_{\boldsymbol{\alpha}} \widetilde{\varphi}_{t,s}(D; \boldsymbol{\alpha}^{1,*}, \boldsymbol{\gamma}_1^*) \right]
&= \E_t\left[ \nabla_{\boldsymbol{\alpha}} \psi(\cdot; \boldsymbol{\alpha}^{1,*}) \right]
- \E_s\left[ w(\boldsymbol{X}; \boldsymbol{\gamma}_1^*) \nabla_{\boldsymbol{\alpha}} \psi(\cdot; \boldsymbol{\alpha}^{1,*}) \right]
\end{aligned}
\]
where $w(\boldsymbol{X}; \boldsymbol{\gamma}) = I(A=1)/e^{\phi(\boldsymbol{X})^T \boldsymbol{\gamma}}$.
This is exactly $\mathbf{0}$ by Lemma \ref{lem:ortho_foc}(i).
Similarly, the gradient w.r.t. $\boldsymbol{\gamma}$ is:
\[
\begin{aligned}
\E\left[ \nabla_{\boldsymbol{\gamma}} \widetilde{\varphi}_{t,s}(D; \boldsymbol{\alpha}^{1,*}, \boldsymbol{\gamma}_1^*) \right]
&= \E_s\left[
    \nabla_{\boldsymbol{\gamma}}w(\boldsymbol{X}; \boldsymbol{\gamma}_1^*)
    \{Y - \psi(\phi(\boldsymbol{X}); \boldsymbol{\alpha}^{1,*})\}
\right] \\
&= -\E_s\left[
    w(\boldsymbol{X}; \boldsymbol{\gamma}_1^*)
    \phi(\boldsymbol{X})
    \{Y - \psi(\phi(\boldsymbol{X}); \boldsymbol{\alpha}^{1,*})\}
\right]
\end{aligned}
\]
This is exactly $\mathbf{0}$ by Lemma \ref{lem:ortho_foc}(ii).
This proves part 1.

\textbf{Part 2 (Variance): explicit $L_2$-Lipschitz bound.}
Write the difference as
\[
\widetilde{\varphi}_{t,s}(D;\vartheta_n)-\widetilde{\varphi}_{t,s}(D;\vartheta^*)
=
T_t(D)+T_s(D),
\]
where
\[
T_t(D)\coloneqq \frac{I(R=t)}{\rho_t}\left\{\psi(\phi(\boldsymbol{X});\boldsymbol{\alpha}_n^1)-\psi(\phi(\boldsymbol{X});\boldsymbol{\alpha}^{1,*})\right\},
\]
and
\[
\begin{aligned}
T_s(D)
&\coloneqq \frac{I(R=s)}{\rho_s}\Big(
w(\boldsymbol{X};\boldsymbol{\gamma}_{1,n})\{Y-\psi(\phi(\boldsymbol{X});\boldsymbol{\alpha}_n^1)\}
-w(\boldsymbol{X};\boldsymbol{\gamma}_1^*)\{Y-\psi(\phi(\boldsymbol{X});\boldsymbol{\alpha}^{1,*})\}
\Big) \\
&= \frac{I(R=s)}{\rho_s}\Big(
\underbrace{\{w(\boldsymbol{X};\boldsymbol{\gamma}_{1,n})-w(\boldsymbol{X};\boldsymbol{\gamma}_1^*)\}\{Y-\psi(\phi(\boldsymbol{X});\boldsymbol{\alpha}^{1,*})\}}_{T_{s,1}(D)}
\;+\;
\underbrace{w(\boldsymbol{X};\boldsymbol{\gamma}_{1,n})\{\psi(\phi(\boldsymbol{X});\boldsymbol{\alpha}^{1,*})-\psi(\phi(\boldsymbol{X});\boldsymbol{\alpha}_n^1)\}}_{T_{s,2}(D)}
\Big).
\end{aligned}
\]
By $(a+b+c)^2\le 3(a^2+b^2+c^2)$ and $\Var(Z)\le \E[Z^2]$,
\[
\Var\!\left(\widetilde{\varphi}_{t,s}(D;\vartheta_n)-\widetilde{\varphi}_{t,s}(D;\vartheta^*)\right)
\le
3\E[T_t(D)^2]+3\E[T_{s,1}(D)^2]+3\E[T_{s,2}(D)^2].
\]
For $T_t$ and $T_{s,2}$, the bounded-slope assumption on $\psi$ (Assumption \ref{asmp:overlap}(b)(i)) implies
$|\psi(u_1)-\psi(u_2)|\le L_\psi|u_1-u_2|$, hence
\[
\E[T_t(D)^2]+\E[T_{s,2}(D)^2]
\le
C\,\E\!\left[\|\phi(\boldsymbol{X})\|_2^2\right]\|\Delta_\alpha\|_2^2,
\]
using the weight-moment condition (Assumption \ref{asmp:overlap}(a)(ii)) and finite second moments of $\phi(\boldsymbol{X})$ (Assumption \ref{asmp:data}(a)).
For $T_{s,1}$, local $L_2$-Lipschitz continuity of $w(\boldsymbol{X};\boldsymbol{\gamma})$ in $\boldsymbol{\gamma}$ (implied by the exponential form of $w$, the overlap/weight-moment assumptions, and the moment condition in Assumption \ref{asmp:overlap}(b)) yields
\[
\E[T_{s,1}(D)^2]
\le
C\|\Delta_\gamma\|_2^2.
\]
Combining the bounds gives the stated variance inequality.
\end{proof}


\begin{lemma}[Double Robustness of Bias Terms]\label{lem:dr_terms_revised}
Let $\Delta_\alpha = \boldsymbol{\alpha}_n^1 - \boldsymbol{\alpha}^{1,*}$ and $\Delta_\gamma = \boldsymbol{\gamma}_{1,n} - \boldsymbol{\gamma}_1^*$. Define the two first-order bias components:
\begin{align*}
\Psi_1(\boldsymbol{\alpha}_n^1, \boldsymbol{\gamma}_1^*) &\coloneqq \E[\widetilde{\varphi}_{t,s}(D; \boldsymbol{\alpha}_n^1, \boldsymbol{\gamma}_1^*) - \widetilde{\varphi}_{t,s}(D; \boldsymbol{\alpha}^{1,*}, \boldsymbol{\gamma}_1^*)] \\
\Psi_2(\boldsymbol{\gamma}_{1,n}, \boldsymbol{\alpha}^{1,*}) &\coloneqq \E[\widetilde{\varphi}_{t,s}(D; \boldsymbol{\alpha}^{1,*}, \boldsymbol{\gamma}_{1,n}) - \widetilde{\varphi}_{t,s}(D; \boldsymbol{\alpha}^{1,*}, \boldsymbol{\gamma}_1^*)]
\end{align*}
Under the conditions:
\begin{enumerate}
    \item If Lemma \ref{lem:ortho_foc}(i) holds (from calibration w.r.t.\ $\boldsymbol{\gamma}$), then $\Psi_1(\boldsymbol{\alpha}_n^1, \boldsymbol{\gamma}_1^*) = O(\lVert \Delta_\alpha \rVert_2^2)$.
    \item If Lemma \ref{lem:ortho_foc}(ii) holds (from calibration w.r.t.\ $\boldsymbol{\alpha}$), then $\Psi_2(\boldsymbol{\gamma}_{1,n}, \boldsymbol{\alpha}^{1,*}) = O(\lVert \Delta_\gamma \rVert_2^2)$.
\end{enumerate}
Note: For identification of the target parameter, Theorem \ref{thm:main} still requires at least one nuisance model to be correctly specified so that $\mu^1(\vartheta^*)=\mu^1$.
This is formalized in Lemma \ref{lem:ident_dr_mu}.
\end{lemma}

\begin{proof}

\textbf{1. Analyzing $\Psi_1$ using Lemma \ref{lem:ortho_foc}(i) (mean-value expansion).}
Let $h_1(t)\coloneqq \E[\widetilde{\varphi}_{t,s}(D;\boldsymbol{\alpha}^{1,*}+t\Delta_\alpha,\boldsymbol{\gamma}_1^*)]$ for $t\in[0,1]$.
Then $\Psi_1=h_1(1)-h_1(0)$ and
\[
\Psi_1
=
\int_0^1 h_1'(t)\,dt
=
\int_0^1 \E\!\left[\nabla_{\boldsymbol{\alpha}}\widetilde{\varphi}_{t,s}(D;\boldsymbol{\alpha}^{1,*}+t\Delta_\alpha,\boldsymbol{\gamma}_1^*)\right]^T\Delta_\alpha\,dt.
\]
Add and subtract $\E[\nabla_{\boldsymbol{\alpha}}\widetilde{\varphi}_{t,s}(D;\boldsymbol{\alpha}^{1,*},\boldsymbol{\gamma}_1^*)]$ inside the integral to get
\[
\Psi_1
=
\left(\E\!\left[\nabla_{\boldsymbol{\alpha}}\widetilde{\varphi}_{t,s}(D;\boldsymbol{\alpha}^{1,*},\boldsymbol{\gamma}_1^*)\right]\right)^T\Delta_\alpha
\;+\;
\int_0^1\left\{\E\!\left[\nabla_{\boldsymbol{\alpha}}\widetilde{\varphi}_{t,s}(D;\boldsymbol{\alpha}^{1,*}+t\Delta_\alpha,\boldsymbol{\gamma}_1^*)\right]
-\E\!\left[\nabla_{\boldsymbol{\alpha}}\widetilde{\varphi}_{t,s}(D;\boldsymbol{\alpha}^{1,*},\boldsymbol{\gamma}_1^*)\right]\right\}^T\Delta_\alpha\,dt.
\]
The first term is zero by Lemma \ref{lem:ortho_foc}(i). For the second term, we bound the gradient difference.
From the definition of $\widetilde{\varphi}_{t,s}$ and the canonical-link GLM identity
$\nabla_{\boldsymbol{\alpha}}\psi(\phi(\boldsymbol{X});\boldsymbol{\alpha})=\phi(\boldsymbol{X})\psi'(\phi(\boldsymbol{X})^T\boldsymbol{\alpha})$,
\[
\nabla_{\boldsymbol{\alpha}}\widetilde{\varphi}_{t,s}(D;\boldsymbol{\alpha},\boldsymbol{\gamma}_1^*)
=
\frac{I(R=t)}{\rho_t}\phi(\boldsymbol{X})\psi'(\phi(\boldsymbol{X})^T\boldsymbol{\alpha})
-\frac{I(R=s)}{\rho_s}w(\boldsymbol{X};\boldsymbol{\gamma}_1^*)\phi(\boldsymbol{X})\psi'(\phi(\boldsymbol{X})^T\boldsymbol{\alpha}).
\]
Using the Lipschitz property of $\psi'$ (Assumption \ref{asmp:overlap}(b)(i)) and the weight-moment condition (Assumption \ref{asmp:overlap}(a)(ii)), we have
\[
\left\|
\E\!\left[\nabla_{\boldsymbol{\alpha}}\widetilde{\varphi}_{t,s}(D;\boldsymbol{\alpha}^{1,*}+t\Delta_\alpha,\boldsymbol{\gamma}_1^*)\right]
-\E\!\left[\nabla_{\boldsymbol{\alpha}}\widetilde{\varphi}_{t,s}(D;\boldsymbol{\alpha}^{1,*},\boldsymbol{\gamma}_1^*)\right]
\right\|_2
\le
C\,\|\Delta_\alpha\|_2,
\]
uniformly over $t\in[0,1]$, where $C<\infty$ depends on $\rho_t,\rho_s$, the moment bound on $\phi(\boldsymbol{X})$ (Assumption \ref{asmp:data}(a)), and $L_{\psi'}$.
Therefore, by Cauchy--Schwarz,
\[
|\Psi_1|
\le
\int_0^1 C\|\Delta_\alpha\|_2 \cdot \|\Delta_\alpha\|_2\,dt
=
O(\|\Delta_\alpha\|_2^2).
\]

\textbf{2. Analyzing $\Psi_2$ using Lemma \ref{lem:ortho_foc}(ii) (mean-value expansion).}
Let $h_2(t)\coloneqq \E[\widetilde{\varphi}_{t,s}(D;\boldsymbol{\alpha}^{1,*},\boldsymbol{\gamma}_1^*+t\Delta_\gamma)]$.
Then $\Psi_2=h_2(1)-h_2(0)$ and
\[
\Psi_2
=
\int_0^1 \E\!\left[\nabla_{\boldsymbol{\gamma}}\widetilde{\varphi}_{t,s}(D;\boldsymbol{\alpha}^{1,*},\boldsymbol{\gamma}_1^*+t\Delta_\gamma)\right]^T\Delta_\gamma\,dt.
\]
Add and subtract the gradient at $\boldsymbol{\gamma}_1^*$:
\[
\Psi_2
=
\left(\E\!\left[\nabla_{\boldsymbol{\gamma}}\widetilde{\varphi}_{t,s}(D;\boldsymbol{\alpha}^{1,*},\boldsymbol{\gamma}_1^*)\right]\right)^T\Delta_\gamma
\;+\;
\int_0^1\left\{\E\!\left[\nabla_{\boldsymbol{\gamma}}\widetilde{\varphi}_{t,s}(D;\boldsymbol{\alpha}^{1,*},\boldsymbol{\gamma}_1^*+t\Delta_\gamma)\right]
-\E\!\left[\nabla_{\boldsymbol{\gamma}}\widetilde{\varphi}_{t,s}(D;\boldsymbol{\alpha}^{1,*},\boldsymbol{\gamma}_1^*)\right]\right\}^T\Delta_\gamma\,dt.
\]
The first term is zero by Lemma \ref{lem:ortho_foc}(ii). For the remainder, recall
$w(\boldsymbol{X};\boldsymbol{\gamma})=I(A=1)\exp(-\phi(\boldsymbol{X})^T\boldsymbol{\gamma})$ so that $\nabla_{\boldsymbol{\gamma}}w(\boldsymbol{X};\boldsymbol{\gamma})=-w(\boldsymbol{X};\boldsymbol{\gamma})\phi(\boldsymbol{X})$.
Thus
\[
\nabla_{\boldsymbol{\gamma}}\widetilde{\varphi}_{t,s}(D;\boldsymbol{\alpha}^{1,*},\boldsymbol{\gamma})
=
\frac{I(R=s)}{\rho_s}\nabla_{\boldsymbol{\gamma}}w(\boldsymbol{X};\boldsymbol{\gamma})\{Y-\psi(\phi(\boldsymbol{X});\boldsymbol{\alpha}^{1,*})\}.
\]
Under Assumption \ref{asmp:overlap}(a,b), the mapping $\boldsymbol{\gamma}\mapsto \E[\nabla_{\boldsymbol{\gamma}}\widetilde{\varphi}_{t,s}(D;\boldsymbol{\alpha}^{1,*},\boldsymbol{\gamma})]$ is locally Lipschitz in $\boldsymbol{\gamma}$, hence
\[
\left\|
\E\!\left[\nabla_{\boldsymbol{\gamma}}\widetilde{\varphi}_{t,s}(D;\boldsymbol{\alpha}^{1,*},\boldsymbol{\gamma}_1^*+t\Delta_\gamma)\right]
-\E\!\left[\nabla_{\boldsymbol{\gamma}}\widetilde{\varphi}_{t,s}(D;\boldsymbol{\alpha}^{1,*},\boldsymbol{\gamma}_1^*)\right]
\right\|_2
\le
C\|\Delta_\gamma\|_2
\]
uniformly over $t\in[0,1]$. Therefore $|\Psi_2|\le \int_0^1 C\|\Delta_\gamma\|_2^2\,dt = O(\|\Delta_\gamma\|_2^2)$.

\textbf{Clarification on Double Robustness:} The key insight is that the calibrated loss functions are designed such that:
\begin{itemize}
    \item Lemma \ref{lem:ortho_foc}(i) holds \textit{by construction} of $L_\gamma$ at any $\vartheta^*$ that minimizes $L_\gamma$, regardless of whether the weighting model is correct.
    \item Lemma \ref{lem:ortho_foc}(ii) holds \textit{by construction} of $L_\alpha$ at any $\vartheta^*$ that minimizes $L_\alpha$, regardless of whether the outcome model is correct.
\end{itemize}
Thus the first-order sensitivity to each nuisance can be removed by calibration; combined with model correctness of at least one nuisance for identification, this yields the double-robust conclusion in Theorem \ref{thm:main}.
\end{proof}

%%%%%%%%%%%%%%%%%%%%%%%%%%%%%%%%%%%%%%%%%%%%%%%%%%%%%%%%%%%%%%%%%%%%%%%%%%%%%%%
%% SECTION 2: AGGREGATED ESTIMATOR PROOFS
%%%%%%%%%%%%%%%%%%%%%%%%%%%%%%%%%%%%%%%%%%%%%%%%%%%%%%%%%%%%%%%%%%%%%%%%%%%%%%%
\section{Aggregated Estimator: FACE-HD}
\label{sec:agg_proofs}

\subsection{Setup and Notation for Aggregation}

Following Section 4.3 in the main manuscript, we aggregate the source-assisted estimators $\{\widehat{\mu}^1_{t,s_j}\}_{j=1}^K$ with the target-only estimator $\widehat{\mu}^1_{ot}$ using an adaptive weighting scheme. Let $K$ be the number of source sites.
Define the global sample size across all participating sites as
\[
N_{\mathrm{all}} := N_t + \sum_{j=1}^K N_{s_j},
\]
and assume $N_t/N_{\mathrm{all}} \to \rho_t^{\mathrm{all}} \in (0,1)$ and $N_{s_j}/N_{\mathrm{all}} \to \rho_{s_j}^{\mathrm{all}} \in (0,1)$ for each fixed $j$.

\textbf{Notation Correspondence with Main Manuscript:}
\begin{itemize}
    \item $\widehat{\mu}^1_{ot}$: Target-only estimator for $\mu^1 = \E_t[Y(1)]$ (uses only target site data)
    \item $\widehat{\mu}^1_{t,s_j}$: Source-assisted estimator using source site $s_j$ (from Section~\ref{sec:pairwise})
    \item $\widehat{\mu}^1_{\text{agg}}$: Aggregated FACE-HD estimator
    \item $\eta_j$: Weight for source site $s_j$
    \item $\varphi_{ot,i}$: Population IF contribution for the target-only estimator (population analogue of $\widehat{\varphi}_{ot,i}$ in the main manuscript)
\end{itemize}

The FACE-HD aggregated estimator (Equation~\eqref{eq:final_opt} in main manuscript) is:
\[
\widehat{\mu}^1_{\text{agg}} = \left(1 - \sum_{j=1}^K \eta_j\right) \widehat{\mu}^1_{ot} + \sum_{j=1}^K \eta_j \widehat{\mu}^1_{t,s_j}
\]
which can be rewritten as:
\[
\widehat{\mu}^1_{\text{agg}} = \widehat{\mu}^1_{ot} + \sum_{j=1}^K \eta_j \left( \widehat{\mu}^1_{t,s_j} - \widehat{\mu}^1_{ot} \right)
\]
where $\widehat{\eta} = (\widehat{\eta}_1, \ldots, \widehat{\eta}_K)^T$ are adaptive weights minimizing the estimated asymptotic variance with a selection penalty.

\subsection{Structure of the Aggregated Asymptotic Covariance}\label{sec:agg_cov_structure}

Before turning to the penalized estimation of $\boldsymbol{\eta}$, we record a structural result on the population covariance entering the aggregation objective. The result clarifies why cross-site covariance terms only couple through the \emph{shared target sample} and justifies the block decomposition used implicitly in Equation~\eqref{eq:final_opt} of the main manuscript.

For each source $s_j$, decompose the pairwise influence-function contribution into an outcome-regression part and a weighted-residual part:
\[
\varphi_{t,s_j}(D;\vartheta^*_{s_j})
= \underbrace{\frac{I(R=t)}{\rho_t}\bigl\{\psi(\phi(\boldsymbol{X});\boldsymbol{\alpha}^{1,*}_{t,s_j}) - \mu^1_t(\boldsymbol{\alpha}^{1,*}_{t,s_j})\bigr\}}_{\displaystyle =:\ \xi^{\mathrm{OR}}_{t,s_j}(D)}
+ \underbrace{\frac{I(R=s_j)}{\rho_{s_j}}\bigl\{w(\boldsymbol{X};\boldsymbol{\gamma}_{s_j,1}^*)(Y-\psi(\phi(\boldsymbol{X});\boldsymbol{\alpha}^{1,*}_{t,s_j})) - \delta^1_{s_j}(\vartheta^*_{s_j})\bigr\}}_{\displaystyle =:\ \xi^{\mathrm{IPW}}_{s_j}(D)}.
\]
For the target-only estimator, denote analogously by $\varphi_{ot}(D)$ the AIPW influence function built from the target-site propensity $\pi_{ot}$ and a target-only outcome model, and by $\xi^{\mathrm{OR}}_{ot}(D)$ and $\xi^{\mathrm{IPW}}_{ot}(D)$ its outcome and residual components.

\begin{lemma}[Block structure of the asymptotic covariance]\label{lem:cov_block_structure}
Suppose the pairwise estimators $\widehat{\mu}^1_{ot}$ and $\{\widehat{\mu}^1_{t,s_j}\}_{j=1}^K$ admit the influence-function representations given above, with $\mathbb{E}[\xi^{\mathrm{OR}}_{\cdot}(D)]=0$ and $\mathbb{E}[\xi^{\mathrm{IPW}}_{\cdot}(D)\mid R=s_j]=0$ (the latter by the calibration moment condition). Then, for any $j,k\in\{1,\ldots,K\}$ with $j\ne k$:
\begin{enumerate}[label={(\alph*)},leftmargin=*]
\item \textbf{(Source residuals are mutually uncorrelated.)} Because observations with $R=s_j$ and $R=s_k$ are disjoint,
\(
\mathbb{E}\bigl[\xi^{\mathrm{IPW}}_{s_j}(D)\,\xi^{\mathrm{IPW}}_{s_k}(D)\bigr]=0.
\)

\item \textbf{(Source residuals are uncorrelated with every OR component.)} By the indicator $I(R=s_j)$ and the target-site support of $\xi^{\mathrm{OR}}_{t,s_k}$,
\(
\mathbb{E}\bigl[\xi^{\mathrm{IPW}}_{s_j}(D)\,\xi^{\mathrm{OR}}_{t,s_k}(D)\bigr]
=\mathbb{E}\bigl[\xi^{\mathrm{IPW}}_{s_j}(D)\,\xi^{\mathrm{OR}}_{ot}(D)\bigr]=0,
\)
for any $j,k$.

\item \textbf{(Cross-site correlation is driven solely by the shared target sample.)} Consequently,
\[
\Cov\bigl(\varphi_{t,s_j}(D),\ \varphi_{t,s_k}(D)\bigr)
=\Cov\bigl(\xi^{\mathrm{OR}}_{t,s_j}(D),\ \xi^{\mathrm{OR}}_{t,s_k}(D)\bigr),
\quad j\ne k,
\]
and similarly the covariance between $\varphi_{ot}$ and any $\varphi_{t,s_j}$ reduces to the covariance of their OR components evaluated on the target sample.
\end{enumerate}
\end{lemma}

\begin{proof}
Parts (a) and (b) follow from $I(R=s_j)I(R=s_k)=0$ for $j\ne k$, from $I(R=s_j)I(R=t)=0$, and from $\mathbb{E}[\xi^{\mathrm{IPW}}_{s_j}(D)\mid R=s_j]=0$. Part (c) follows by expanding the covariance using the decomposition of each $\varphi_{t,s_j}$ into its OR and IPW components and applying (a)--(b) to cancel the IPW cross-terms.
\end{proof}

The block structure established by Lemma~\ref{lem:cov_block_structure} implies that the population Hessian $H$ in the quadratic aggregation objective $L^*(\boldsymbol{\eta})=\boldsymbol{\eta}^\top H\boldsymbol{\eta}-2\,g^\top\boldsymbol{\eta}+c$ decomposes as $H=H^{\mathrm{OR}}+H^{\mathrm{IPW}}$, where $H^{\mathrm{IPW}}$ is diagonal (containing only the per-source residual variances) and $H^{\mathrm{OR}}$ is the $K\times K$ Gram matrix induced by the shared target-sample OR components. This decomposition is the direct high-dimensional counterpart of Equation~(11) in \citet{han2025federated}: the first sum $\sum_{k_s\in\mathcal{S}}\eta_{k_s}^2\,\widehat{\sigma}^2_{k_s}/n_{k_s}$ in the FACE aggregation objective corresponds to the diagonal $H^{\mathrm{IPW}}$ contribution, and the second sum $\sum_{k_t\in\mathcal{T}}\widehat{h}_{k_t}(\boldsymbol{\eta})^\top\widehat{\Sigma}_{k_t}/n_{k_t}\,\widehat{h}_{k_t}(\boldsymbol{\eta})$ encodes the cross-source coupling through the shared target sample that corresponds to $H^{\mathrm{OR}}$. The full penalized QP in Equation~\eqref{eq:final_opt} of the main manuscript, like the FACE objective, does not admit a closed-form minimizer in general and is solved numerically. The following corollary records the explicit solution in the limiting regime where the target-sample coupling vanishes, which is the high-dimensional analogue of the same diagonal simplification implicit in the coordinate-wise structure of the FACE aggregation rule.

\begin{corollary}[Soft-thresholding closed form under diagonal Hessian]\label{cor:soft_threshold_face}
Let $\Delta_j\coloneqq \widehat{\mu}^1_{t,s_j}-\widehat{\mu}^1_{ot}$ and $d_j\coloneqq \widehat{\mu}^1_{ot}-\widehat{\mu}^1_{t,s_j}=-\Delta_j$. Suppose that, in addition to the hypotheses of Lemma~\ref{lem:cov_block_structure}, the OR components are pairwise uncorrelated across source sites in the sense that
\begin{equation}\label{eq:diag_cond}
\Cov\!\bigl(\xi^{\mathrm{OR}}_{t,s_j}(D),\ \xi^{\mathrm{OR}}_{t,s_k}(D)\bigr)=0\qquad\text{for all }j\neq k.
\end{equation}
Then the population Hessian $H$ is diagonal with entries
\[
H_{jj}\;=\;\Var(\Delta_j)\;=\;\sigma_j^2,\qquad g_j\;=\;-\Cov(\widehat{\mu}^1_{ot},\,\Delta_j),
\]
where $\sigma_j^2$ collects the pairwise and target-only variance components present in Equation~\eqref{eq:final_opt} of the main manuscript, and the penalized aggregation problem
\[
\widehat{\boldsymbol{\eta}}\;=\;\arg\min_{\boldsymbol{\eta}\in\mathbb{R}^K}\left\{\boldsymbol{\eta}^\top\widehat{H}\boldsymbol{\eta}-2\widehat{g}^\top\boldsymbol{\eta}+\lambda\sum_{j=1}^K|\eta_j|\,d_j^2\right\}
\]
with $\widehat{H}$ diagonal admits the coordinate-wise closed-form \emph{soft-thresholding} solution
\begin{equation}\label{eq:soft_thresh}
\widehat{\eta}_j
\;=\;
\frac{1}{\widehat{H}_{jj}}\,\operatorname{sign}(\widehat{g}_j)\,\max\!\left\{0,\ |\widehat{g}_j|-\tfrac{\lambda}{2}\,d_j^2\right\}
\;=\;
\operatorname{sign}(\widehat{\eta}_j^{\mathrm{OLS}})\,\max\!\left\{0,\ |\widehat{\eta}_j^{\mathrm{OLS}}|-\frac{\lambda\,d_j^2}{2\,\widehat{H}_{jj}}\right\},
\end{equation}
where $\widehat{\eta}_j^{\mathrm{OLS}}\coloneqq \widehat{g}_j/\widehat{H}_{jj}$ is the unpenalized coordinate minimizer. In particular, $\widehat{\eta}_j=0$ whenever $|\widehat{\eta}_j^{\mathrm{OLS}}|\le \lambda d_j^2/(2\widehat{H}_{jj})$, so the selection step of the aggregation becomes an explicit thresholding rule.
\end{corollary}

\begin{proof}
Under \eqref{eq:diag_cond}, Lemma~\ref{lem:cov_block_structure}(c) gives $\Cov(\Delta_j,\Delta_k)=\Cov(\xi^{\mathrm{OR}}_{t,s_j},\xi^{\mathrm{OR}}_{t,s_k})+\Cov(\xi^{\mathrm{IPW}}_{s_j},\xi^{\mathrm{IPW}}_{s_k})=0$ for $j\neq k$, where the IPW cross-term vanishes by disjoint support. Hence the population (and plug-in empirical) Hessian is diagonal, $\widehat{H}=\mathrm{diag}(\widehat{H}_{11},\ldots,\widehat{H}_{KK})$.

Because $\widehat{H}$ is diagonal and the penalty decouples across coordinates, the aggregation problem separates into $K$ univariate problems:
\begin{equation}\label{eq:coord_obj}
\widehat{\eta}_j\;=\;\arg\min_{\eta\in\mathbb{R}}\,\bigl\{\widehat{H}_{jj}\,\eta^2-2\widehat{g}_j\,\eta+\lambda\,d_j^2\,|\eta|\bigr\}.
\end{equation}
Fix $j$ and let $h\coloneqq\widehat{H}_{jj}>0$, $g\coloneqq\widehat{g}_j$, $\ell\coloneqq\lambda d_j^2/2\ge 0$. The objective in \eqref{eq:coord_obj} is strictly convex, so its minimizer is characterized by the subdifferential optimality condition
\[
0\in\partial\bigl(h\eta^2-2g\eta+2\ell|\eta|\bigr)
=\{2h\eta-2g+2\ell\,v:v\in\operatorname{sgn}(\eta)\},
\]
where $\operatorname{sgn}(\eta)=\{1\}$ if $\eta>0$, $\{-1\}$ if $\eta<0$, and $[-1,1]$ if $\eta=0$.

\emph{Case $\eta>0$.} Optimality requires $h\eta-g+\ell=0$, i.e., $\eta=(g-\ell)/h$. This is consistent with $\eta>0$ iff $g>\ell$, in which case $\eta=(g-\ell)/h>0$ minimizes the objective (by strict convexity).

\emph{Case $\eta<0$.} Optimality requires $h\eta-g-\ell=0$, i.e., $\eta=(g+\ell)/h$. This is consistent with $\eta<0$ iff $g<-\ell$, in which case $\eta=(g+\ell)/h<0$.

\emph{Case $\eta=0$.} Optimality requires $0\in 2h\cdot 0-2g+2\ell\cdot[-1,1]$, i.e., $g\in[-\ell,\ell]$, which is consistent iff $|g|\le\ell$.

Combining the three cases,
\[
\widehat{\eta}_j=\begin{cases}
(g-\ell)/h,&g>\ell,\\
0,&|g|\le\ell,\\
(g+\ell)/h,&g<-\ell,
\end{cases}
\]
which collapses to the single expression $\widehat{\eta}_j=h^{-1}\operatorname{sign}(g)\max\{0,|g|-\ell\}$. Restoring $h=\widehat{H}_{jj}$, $g=\widehat{g}_j$, $\ell=\lambda d_j^2/2$ yields the first form of \eqref{eq:soft_thresh}. The second form is obtained by factoring out $\widehat{H}_{jj}$ using $\widehat{\eta}_j^{\mathrm{OLS}}=\widehat{g}_j/\widehat{H}_{jj}$ and noting that $\operatorname{sign}$ is invariant under multiplication by the positive scalar $\widehat{H}_{jj}$.

Finally, restating the minimizer in terms of the inverse-variance unpenalized weight $\widehat{\eta}_j^{\mathrm{OLS}}$ and the data-adaptive threshold $\lambda d_j^2/(2\widehat{H}_{jj})$ emphasizes the soft-thresholding interpretation: sources whose unpenalized weight is dominated by the squared target--source gap are exactly screened out. This is the high-dimensional counterpart of the selection behavior observed in the coordinate-wise structure of the FACE aggregation rule of \citet{han2025federated}, which, however, retains a non-trivial cross-source coupling through $\widehat{h}_{k_t}(\boldsymbol{\eta})$ (the analogue of $H^{\mathrm{OR}}$) and therefore does not admit a closed form in general.
\end{proof}

\begin{remark}[When does condition \eqref{eq:diag_cond} hold?]\label{rem:diag_condition}
Condition \eqref{eq:diag_cond} is not generic: because $\xi^{\mathrm{OR}}_{t,s_j}(D)$ is supported on the shared target sample, its covariance with $\xi^{\mathrm{OR}}_{t,s_k}(D)$ depends on the correlation structure of the outcome models fitted with different pairwise nuisance targets $\boldsymbol{\alpha}^{1,*}_{t,s_j}$ and $\boldsymbol{\alpha}^{1,*}_{t,s_k}$. The condition holds in two natural limiting regimes: (i) when the per-source residual variances dominate the shared-target cross-covariances so that the off-diagonal entries of $H^{\mathrm{OR}}$ are negligible compared to the diagonal entries of $H^{\mathrm{IPW}}$; (ii) when the pairwise OR targets coincide asymptotically, in which case the shared-target covariance collapses into a rank-one structure that can be absorbed into the intercept. Outside these regimes, Corollary~\ref{cor:soft_threshold_face} provides a first-order approximation whose accuracy depends on the magnitude of the off-diagonal entries of $H^{\mathrm{OR}}$; the general solution, both in FACE and in FACE-HD, requires solving the full $K$-dimensional penalized QP as in Equation~\eqref{eq:final_opt}.
\end{remark}

\subsection{Penalized Variance Estimation}

Define the empirical loss function as the estimated asymptotic variance of the aggregated estimator:
\[
\widehat{L}(\eta) = N_{\mathrm{all}} \cdot \widehat{\Var}\left( \widehat{\mu}^1_{\text{agg}}(\eta) \right)
\]
which is a quadratic function of $\eta$ (equivalently, $N_{\mathrm{all}}$ times the variance expression in Equation~\eqref{eq:final_opt} in the main manuscript):
\[
\widehat{L}(\eta) = \eta^T \widehat{H} \eta + \widehat{g}^T \eta + \widehat{c}
\]
where $\widehat{H}$, $\widehat{g}$, and $\widehat{c}$ are the estimated Hessian, gradient, and constant terms derived from the variance components in Section 4.2 of the main manuscript.

\begin{lemma}[Consistency of Variance/Covariance Components]\label{lem:var_components_agg}
Assume Assumptions \ref{asmp:data} and \ref{asmp:overlap}, and that $K$ is fixed. Let $H,g,c$ denote the population analogues of $\widehat{H},\widehat{g},\widehat{c}$ defined by replacing empirical averages with expectations and using the population influence functions (with the same cross-site covariance terms induced by the shared target sample). Then:
\[
\|\widehat{H}-H\|_{\mathrm{op}} = O_p(N_{\mathrm{all}}^{-1/2}),\qquad
\|\widehat{g}-g\|_2 = O_p(N_{\mathrm{all}}^{-1/2}),\qquad
|\widehat{c}-c| = O_p(N_{\mathrm{all}}^{-1/2}),
\]
and consequently, for any fixed compact set $\mathcal{B}\subset\R^K$,
\[
\sup_{\eta\in\mathcal{B}}|\widehat{L}(\eta)-L^*(\eta)| = O_p(N_{\mathrm{all}}^{-1/2}).
\]
\end{lemma}
\begin{proof}
We prove entrywise $\sqrt{N_{\mathrm{all}}}$-consistency and then convert to the stated matrix/vector norms using that $K$ is fixed.

\textbf{Step 1 (Generic moment form).}
Each element of $\widehat{H}$, $\widehat{g}$ and $\widehat{c}$ is a (finite) linear combination of empirical averages of the form
\[
\widehat{M}_{r} \;=\; \frac{1}{N_r}\sum_{i\in\mathcal{D}_r} f(D_i;\widehat{\vartheta}^{(-\mathrm{fold}(i))}),
\qquad r\in\{t,s_1,\ldots,s_K\},
\]
where $f$ is a product of two empirical influence-function components (including the cross-site covariance terms computed on the \emph{same} target sample), and $\widehat{\vartheta}^{(-\mathrm{fold}(i))}$ denotes the cross-fitted nuisance estimate used to evaluate the fold containing $i$.
Let $M_r$ be the population analogue obtained by replacing the sample average with $\E_r[\cdot]$ and $\widehat{\vartheta}^{(-\mathrm{fold}(i))}$ with the population target $\vartheta^*$.

\textbf{Step 2 (Conditional LLN and plug-in stability).}
By cross-fitting, conditional on the nuisance-training samples, the summands $\{f(D_i;\widehat{\vartheta}^{(-\mathrm{fold}(i))}):i\in\mathcal{D}_r\}$ are independent within each site $r$, with conditional mean $\E_r[f(D;\widehat{\vartheta}^{(-\mathrm{fold}(D))})\mid \{\mathcal{D}^{(-k)}\}]$.
Assumption \ref{asmp:overlap}(b) gives a finite $(1+\delta)$ moment for $f(D;\vartheta)$ uniformly in a neighborhood of $\vartheta^*$ (since it is a product of two $(2+\delta)$-moment IF components), hence the conditional variance is finite and a conditional LLN/CLT implies
\[
\widehat{M}_r - \E_r\!\left[f(D;\widehat{\vartheta}^{(-\mathrm{fold}(D))})\,\middle|\, \{\mathcal{D}^{(-k)}\}\right]
= O_p(N_r^{-1/2}).
\]
Next, by Lemma \ref{lem:nuisance_rates}, $\widehat{\vartheta}^{(-k)}\to_p \vartheta^*$, and by Lemma \ref{lem:bias_bound} (Part 2) the mapping $\vartheta\mapsto f(D;\vartheta)$ is locally $L_2$-Lipschitz. Therefore,
\[
\E_r\!\left[f(D;\widehat{\vartheta}^{(-\mathrm{fold}(D))})\right]
=
\E_r[f(D;\vartheta^*)] + o(1),
\]
and combining the two displays yields $\widehat{M}_r-M_r=O_p(N_r^{-1/2})$.
Under the site-ratio condition $N_r/N_{\mathrm{all}}\to\rho_r>0$, this is $O_p(N_{\mathrm{all}}^{-1/2})$.

\textbf{Step 3 (From entrywise to matrix/vector norms).}
Since $\widehat{H}$ has fixed dimension $K\times K$ and $\widehat{g}$ has fixed dimension $K$, the preceding bound applies to each entry, hence
\[
\|\widehat{H}-H\|_{\mathrm{op}}
\le
C_K\max_{j,k}|\widehat{H}_{jk}-H_{jk}|
=
O_p(N_{\mathrm{all}}^{-1/2}),
\qquad
\|\widehat{g}-g\|_2
\le
\sqrt{K}\max_j|\widehat{g}_j-g_j|
=
O_p(N_{\mathrm{all}}^{-1/2}),
\]
Finally, $\widehat{c}$ is a finite linear combination of the same sample averages $\{\widehat{M}_r\}$ (coming from the constant terms in the variance expansion), so the same entrywise $O_p(N_{\mathrm{all}}^{-1/2})$ bound implies $|\widehat{c}-c|=O_p(N_{\mathrm{all}}^{-1/2})$.

\textbf{Step 4 (Uniform convergence of the quadratic form).}
For $\eta$ in a fixed compact set $\mathcal{B}\subset\R^K$,
\[
|\widehat{L}(\eta)-L^*(\eta)|
\le
\|\widehat{H}-H\|_{\mathrm{op}}\|\eta\|_2^2 + \|\widehat{g}-g\|_2\|\eta\|_2 + |\widehat{c}-c|,
\]
and the right-hand side is $O_p(N_{\mathrm{all}}^{-1/2})$ uniformly over $\eta\in\mathcal{B}$ because $\sup_{\eta\in\mathcal{B}}\|\eta\|_2<\infty$.
\end{proof}

The penalized objective from Equation~\eqref{eq:final_opt} is:
\begin{equation}\label{eq:penalized_obj}
\widehat{\eta} = \arg\min_\eta \left\{ \widehat{L}(\eta) + \lambda \sum_{j=1}^K |\eta_j| \left( \widehat{\mu}^1_{ot} - \widehat{\mu}^1_{t,s_j} \right)^2 \right\}
\end{equation}
where $\lambda > 0$ is a tuning parameter that depends on $N_{\mathrm{all}}$.

\textbf{Remark:} The penalty term $(\widehat{\mu}^1_{ot} - \widehat{\mu}^1_{t,s_j})^2$ measures the discrepancy between the target-only and source-assisted estimators. Large discrepancies indicate potential covariate shift or model misspecification, and the penalty drives the corresponding $\eta_j$ toward zero.

\subsection{Informative Source Sites}

Let $\mu^1 = \E_t[Y(1)]$ be the true target-site potential outcome mean.

Define the set of \textit{informative} source sites as:
\[
\mathcal{S}^* = \left\{ j \in \{1, \ldots, K\} : \text{plim}\ \widehat{\mu}^1_{t,s_j} = \mu^1 \right\}
\]
i.e., sites where the source-assisted estimator is consistent for the target parameter $\mu^1$.

Define the oracle-restricted space:
\[
\mathbb{R}_{\mathcal{S}^*} := \left\{ \eta \in \mathbb{R}^K : \eta_j = 0 \text{ for all } j \notin \mathcal{S}^* \right\}.
\]
Let $\eta^*$ be the oracle optimal weights that minimize the population asymptotic variance over $\mathbb{R}_{\mathcal{S}^*}$:
\[
\eta^* = \arg\min_{\eta \in \mathbb{R}_{\mathcal{S}^*}} L^*(\eta), \quad
L^*(\eta) := \Var\left( \sqrt{N_{\mathrm{all}}}(\widehat{\mu}^1_{\text{agg}}(\eta) - \mu^1) \right).
\]

\begin{assumption}[Aggregation Regularity]\label{asmp:aggregation}
(a) (Sufficient Informative Sources + Local Strong Convexity) The set $\mathcal{S}^*$ is non-empty, i.e., there exists at least one informative source site. Moreover, the population aggregation objective is locally strictly convex on the oracle subspace: the Hessian of $L^*(\eta)$ restricted to $\mathbb{R}_{\mathcal{S}^*}$, denoted $H_{\mathcal{S}^*}$, is positive definite (equivalently, $\eta^*$ is locally unique).
(b) (Separation / ``$\beta$-min'' for biased sites) Define
\[
b_{0,N} \coloneqq \min_{j\notin\mathcal{S}^*}\left|\text{plim}(\widehat{\mu}^1_{t,s_j}) - \mu^1\right|.
\]
Assume $\lambda\,b_{0,N}^2\to\infty$. A special case is \emph{fixed separation}: $b_{0,N}\ge b_0>0$.
(c) (Penalty Growth) The penalty parameter satisfies $\lambda \to \infty$ and $\lambda / \sqrt{N_{\mathrm{all}}} \to 0$.
(d) (Fixed $K$) The number of source sites $K$ is fixed.
\end{assumption}

\begin{remark}[When separation fails / local alternatives]\label{rem:weak_bias}
Assumption \ref{asmp:aggregation}(b) is used only for \emph{support recovery} (forcing $\widehat{\eta}_j=0$ for biased sites). If some biased sites are only \emph{weakly} biased in the sense that $b_{0,N}\to 0$, then selection consistency may fail. In that regime, the aggregated estimator remains asymptotically linear for any deterministic $\eta$ (Lemma \ref{lem:if_agg}), and the additional first-order bias is controlled by
\[
\left|\sum_{j\notin\mathcal{S}^*}\eta_j\{\text{plim}(\widehat{\mu}^1_{t,s_j})-\mu^1\}\right|,
\]
so if any included weakly biased sites satisfy $b_{j,N}=o(N_{\mathrm{all}}^{-1/2})$, their contribution is still negligible at $\sqrt{N_{\mathrm{all}}}$ scale.
\end{remark}

\subsection{Technical Lemmas for Aggregation}

\begin{lemma}[Influence Function Representation]\label{lem:if_agg}
Under Assumptions \ref{asmp:data}-\ref{asmp:rsc} and \ref{asmp:aggregation}(a), the aggregated estimator admits the asymptotic linear representation:
\[
\sqrt{N_{\mathrm{all}}}\left( \widehat{\mu}^1_{\text{agg}}(\eta^*) - \mu^1 \right) = \frac{1}{\sqrt{N_{\mathrm{all}}}} \sum_{i=1}^{N_{\mathrm{all}}} \varphi_{\mathrm{agg},i}(\eta^*) + o_p(1)
\]
where $\varphi_{\mathrm{agg},i}(\eta^*)$ is the aggregated influence function:
\[
\varphi_{\mathrm{agg},i}(\eta^*) = \left(1-\sum_{j\in\mathcal{S}^*}\eta_j^*\right)\varphi_{ot,i} + \sum_{j\in\mathcal{S}^*}\eta_j^* \varphi_{t,s_j,i},
\]
with $\varphi_{t,s_j,i}$ denoting the (rescaled) IF contribution of source-assisted estimator $j$ under $\sqrt{N_{\mathrm{all}}}$ scaling.
\textbf{Convention:} Each $\varphi_{ot,i}$ (resp.\ $\varphi_{t,s_j,i}$) is understood as zero for observations not in the target data $\mathcal{D}_t$ (resp.\ not in $\mathcal{D}_t\cup\mathcal{D}_{s_j}$), and includes the constant rescaling needed to convert its natural $\sqrt{N_t}$ (resp.\ $\sqrt{N_{t,s_j}}$) linear representation into the common $\sqrt{N_{\mathrm{all}}}$ scaling.
\end{lemma}

\begin{proof}
By Theorem \ref{thm:main}, for each fixed $j\in\mathcal{S}^*$, there is an asymptotic linear representation under pairwise scaling $N_{t,s_j}:=N_t+N_{s_j}$:
\[
\sqrt{N_{t,s_j}}(\widehat{\mu}^1_{t,s_j} - \text{plim}(\widehat{\mu}^1_{t,s_j})) = \frac{1}{\sqrt{N_{t,s_j}}} \sum_{i=1}^{N_{\mathrm{all}}} \bar{\varphi}_{t,s_j,i} + o_p(1),
\]
where $\bar{\varphi}_{t,s_j,i}$ is understood as zero for observations not in $\mathcal{D}_t\cup\mathcal{D}_{s_j}$. Since $N_{t,s_j}/N_{\mathrm{all}}\to c_j\in(0,1)$, the above is equivalent to
\[
\sqrt{N_{\mathrm{all}}}(\widehat{\mu}^1_{t,s_j} - \text{plim}(\widehat{\mu}^1_{t,s_j})) = \frac{1}{\sqrt{N_{\mathrm{all}}}} \sum_{i=1}^{N_{\mathrm{all}}} \varphi_{t,s_j,i} + o_p(1)
\]
for a rescaled IF $\varphi_{t,s_j,i}$. The target-only estimator has the analogous form:
\[
\sqrt{N_{\mathrm{all}}}(\widehat{\mu}^1_{ot} - \mu^1) = \frac{1}{\sqrt{N_{\mathrm{all}}}} \sum_{i=1}^{N_{\mathrm{all}}} \varphi_{ot,i} + o_p(1).
\]
For $j \in \mathcal{S}^*$, $\text{plim}(\widehat{\mu}^1_{t,s_j})=\mu^1$ by definition.

By linearity, the aggregated estimator with fixed $\eta = \eta^*$ satisfies:
\[
\sqrt{N_{\mathrm{all}}}(\widehat{\mu}^1_{\text{agg}}(\eta^*) - \mu^1) = \frac{1}{\sqrt{N_{\mathrm{all}}}} \sum_{i=1}^{N_{\mathrm{all}}} \varphi_{\mathrm{agg},i}(\eta^*) + o_p(1).
\]
The remainder remains $o_p(1)$ because $K$ is fixed.
\end{proof}

\begin{lemma}[Selection Consistency]\label{lem:selection}
Under Assumptions \ref{asmp:data}-\ref{asmp:rsc} and \ref{asmp:aggregation}, the penalized estimator $\widehat{\eta}$ from \eqref{eq:penalized_obj} satisfies:
\begin{enumerate}
    \item \textbf{(Rate)} $\|\widehat{\eta} - \eta^*\|_2 = O_p(N_{\mathrm{all}}^{-1/2})$
    \item \textbf{(Selection)} $\widehat{\eta}_j = 0$ for all $j \notin \mathcal{S}^*$ with probability approaching 1.
\end{enumerate}
\end{lemma}

\begin{proof}
We give a self-contained argument; it parallels the FACE proof in \citet{han2025federated} (Lemma S5) but the objective here is an explicit quadratic form.
Recall $\widehat{L}(\eta)=\eta^T\widehat{H}\eta+\widehat{g}^T\eta+\widehat{c}$ and its population analogue $L^*(\eta)=\eta^TH\eta+g^T\eta+c$.

\textbf{Part 1 (Rate on the oracle subspace).}
Define the oracle-restricted estimator
\[
\widetilde{\eta} \;=\; \arg\min_{\eta\in\mathbb{R}_{\mathcal{S}^*}}
\left\{ \widehat{L}(\eta) + \lambda \sum_{j=1}^K |\eta_j| (\widehat{\mu}^1_{ot} - \widehat{\mu}^1_{t,s_j})^2 \right\}.
\]
On $\mathbb{R}_{\mathcal{S}^*}$, only coordinates $j\in\mathcal{S}^*$ can be nonzero, so the penalty term reduces to $\sum_{j\in\mathcal{S}^*}|\eta_j|(\widehat{\mu}^1_{ot}-\widehat{\mu}^1_{t,s_j})^2$.
By Theorem \ref{thm:main} and the site-ratio regularity, for each fixed $j\in\mathcal{S}^*$,
$
\widehat{\mu}^1_{ot}-\widehat{\mu}^1_{t,s_j}=O_p(N_{\mathrm{all}}^{-1/2}),
$
hence $(\widehat{\mu}^1_{ot}-\widehat{\mu}^1_{t,s_j})^2=O_p(N_{\mathrm{all}}^{-1})$ and therefore the penalty contribution on $\mathbb{R}_{\mathcal{S}^*}$ is
$O_p(\lambda/N_{\mathrm{all}})=o_p(N_{\mathrm{all}}^{-1/2})$ under Assumption \ref{asmp:aggregation}(c).

Next, by Lemma \ref{lem:var_components_agg}, the quadratic coefficients satisfy $\|\widehat{H}-H\|_{\mathrm{op}}=O_p(N_{\mathrm{all}}^{-1/2})$ and $\|\widehat{g}-g\|_2=O_p(N_{\mathrm{all}}^{-1/2})$.
Let $H_{\mathcal{S}^*}$ and $\widehat{H}_{\mathcal{S}^*}$ denote the principal submatrices corresponding to $\mathcal{S}^*$, and similarly for $g$.
Because $K$ is fixed, $H_{\mathcal{S}^*}$ is a fixed-dimensional positive definite matrix (local strong convexity of $L^*$ around $\eta^*$), hence $\|H_{\mathcal{S}^*}^{-1}\|_{\mathrm{op}}<\infty$ and $\|\widehat{H}_{\mathcal{S}^*}^{-1}\|_{\mathrm{op}}=O_p(1)$.
More explicitly, by Weyl's inequality and Lemma \ref{lem:var_components_agg}, $\lambda_{\min}(\widehat{H}_{\mathcal{S}^*})\ge \lambda_{\min}(H_{\mathcal{S}^*}) - \|\widehat{H}_{\mathcal{S}^*}-H_{\mathcal{S}^*}\|_{\mathrm{op}}$, which is bounded away from 0 w.p.a.1, so $\widehat{H}_{\mathcal{S}^*}$ is invertible and $\|\widehat{H}_{\mathcal{S}^*}^{-1}\|_{\mathrm{op}}=O_p(1)$.
Ignoring the $o_p(N_{\mathrm{all}}^{-1/2})$ penalty perturbation, the minimizers of $\widehat{L}$ and $L^*$ on $\mathbb{R}_{\mathcal{S}^*}$ satisfy the first-order conditions
$2\widehat{H}_{\mathcal{S}^*}\eta + \widehat{g}_{\mathcal{S}^*}=0$ and $2H_{\mathcal{S}^*}\eta + g_{\mathcal{S}^*}=0$,
so
\[
\eta^*_{\mathcal{S}^*} = -\tfrac{1}{2}H_{\mathcal{S}^*}^{-1}g_{\mathcal{S}^*},
\qquad
\widetilde{\eta}_{\mathcal{S}^*} = -\tfrac{1}{2}\widehat{H}_{\mathcal{S}^*}^{-1}\widehat{g}_{\mathcal{S}^*} + o_p(N_{\mathrm{all}}^{-1/2}).
\]
To make the perturbation step explicit, write
\[
\widetilde{\eta}_{\mathcal{S}^*}-\eta^*_{\mathcal{S}^*}
=
-\tfrac{1}{2}\left(\widehat{H}_{\mathcal{S}^*}^{-1}\widehat{g}_{\mathcal{S}^*}-H_{\mathcal{S}^*}^{-1}g_{\mathcal{S}^*}\right)
+o_p(N_{\mathrm{all}}^{-1/2}).
\]
Then
\[
\widehat{H}_{\mathcal{S}^*}^{-1}\widehat{g}_{\mathcal{S}^*}-H_{\mathcal{S}^*}^{-1}g_{\mathcal{S}^*}
=
\widehat{H}_{\mathcal{S}^*}^{-1}(\widehat{g}_{\mathcal{S}^*}-g_{\mathcal{S}^*})
+(\widehat{H}_{\mathcal{S}^*}^{-1}-H_{\mathcal{S}^*}^{-1})g_{\mathcal{S}^*}.
\]
The first term is $O_p(N_{\mathrm{all}}^{-1/2})$ since $\|\widehat{H}_{\mathcal{S}^*}^{-1}\|_{\mathrm{op}}=O_p(1)$ and $\|\widehat{g}-g\|_2=O_p(N_{\mathrm{all}}^{-1/2})$.
For the second term, use $\widehat{H}_{\mathcal{S}^*}^{-1}-H_{\mathcal{S}^*}^{-1}=\widehat{H}_{\mathcal{S}^*}^{-1}(H_{\mathcal{S}^*}-\widehat{H}_{\mathcal{S}^*})H_{\mathcal{S}^*}^{-1}$, so
\[
\|\widehat{H}_{\mathcal{S}^*}^{-1}-H_{\mathcal{S}^*}^{-1}\|_{\mathrm{op}}
\le
\|\widehat{H}_{\mathcal{S}^*}^{-1}\|_{\mathrm{op}}\|\widehat{H}-H\|_{\mathrm{op}}\|H_{\mathcal{S}^*}^{-1}\|_{\mathrm{op}}
=
O_p(N_{\mathrm{all}}^{-1/2}),
\]
and hence the second term is also $O_p(N_{\mathrm{all}}^{-1/2})$. Therefore $\|\widetilde{\eta}-\eta^*\|_2 = O_p(N_{\mathrm{all}}^{-1/2})$.
Finally, since $\widehat{\eta}$ agrees with $\widetilde{\eta}$ w.p.a.1 by Part 2 below, the same rate holds for $\widehat{\eta}$.

\textbf{Part 2 (Selection of biased sites).}
By Part 1, $\widetilde{\eta}=\eta^*+O_p(N_{\mathrm{all}}^{-1/2})$, so there exists a fixed-radius neighborhood
$\mathcal{B}\coloneqq\{\eta:\|\eta-\eta^*\|_2\le r\}$ such that $P(\widetilde{\eta}\in\mathcal{B})\to 1$.
For the full penalized minimizer, the same neighborhood restriction applies w.p.a.1 once we show (below) that all biased coordinates are set to zero, i.e., $\widehat{\eta}=\widetilde{\eta}$ w.p.a.1. Thus, in the KKT bound we may work on $\mathcal{B}$.

For biased sites $j \notin \mathcal{S}^*$, let $b_{j,N}\coloneqq |\text{plim}(\widehat{\mu}^1_{t,s_j}) - \mu^1|$, so $b_{j,N}\ge b_{0,N}$.
Because $\widehat{\mu}^1_{t,s_j}\xrightarrow{p}\text{plim}(\widehat{\mu}^1_{t,s_j})$ and $\widehat{\mu}^1_{ot}\xrightarrow{p}\mu^1$,
\[
(\widehat{\mu}^1_{ot} - \widehat{\mu}^1_{t,s_j})^2 \xrightarrow{p} b_{j,N}^2 \ge b_{0,N}^2,
\]
and thus $\lambda(\widehat{\mu}^1_{ot} - \widehat{\mu}^1_{t,s_j})^2\to\infty$ under Assumption \ref{asmp:aggregation}(b).

The KKT condition for the penalized objective at coordinate $j$ implies that if $\widehat{\eta}_j=0$ then
\[
\left| \frac{\partial \widehat{L}}{\partial \eta_j}(\widehat{\eta}) \right|
\le \lambda (\widehat{\mu}^1_{ot} - \widehat{\mu}^1_{t,s_j})^2.
\]
Since $\widehat{L}(\eta)=\eta^T\widehat{H}\eta+\widehat{g}^T\eta+\widehat{c}$ is quadratic,
\[
\frac{\partial \widehat{L}}{\partial \eta_j}(\eta)
=
2(\widehat{H}\eta)_j+\widehat{g}_j,
\]
so for any compact $\mathcal{B}\subset\mathbb{R}^K$,
\[
\sup_{\eta\in\mathcal{B}}\left|\frac{\partial \widehat{L}}{\partial \eta_j}(\eta)\right|
\le
2\sup_{\eta\in\mathcal{B}}\|\widehat{H}\eta\|_\infty+\|\widehat{g}\|_\infty
\le
2\|\widehat{H}\|_{\mathrm{op}}\sup_{\eta\in\mathcal{B}}\|\eta\|_2+\|\widehat{g}\|_2.
\]
By Lemma \ref{lem:var_components_agg}, $\|\widehat{H}\|_{\mathrm{op}}=O_p(1)$ and $\|\widehat{g}\|_2=O_p(1)$, hence $\sup_{\eta\in\mathcal{B}}|\partial \widehat{L}/\partial \eta_j(\eta)|=O_p(1)$.
Because the right-hand side diverges, the KKT inequality is satisfied for $\eta_j=0$ w.p.a.1 and any minimizer must set $\widehat{\eta}_j=0$ for all $j\notin\mathcal{S}^*$.
Therefore $\widehat{\eta}\in\mathbb{R}_{\mathcal{S}^*}$ w.p.a.1, hence $\widehat{\eta}=\widetilde{\eta}$ w.p.a.1.
\end{proof}

\begin{lemma}[Asymptotic Equivalence]\label{lem:asymp_equiv}
Under Assumptions \ref{asmp:data}-\ref{asmp:rsc} and \ref{asmp:aggregation}:
\[
\sqrt{N_{\mathrm{all}}}\left( \widehat{\mu}^1_{\text{agg}}(\widehat{\eta}) - \widehat{\mu}^1_{\text{agg}}(\eta^*) \right) = o_p(1)
\]
\end{lemma}

\begin{proof}
The difference is:
\[
\sqrt{N_{\mathrm{all}}}\left( \widehat{\mu}^1_{\text{agg}}(\widehat{\eta}) - \widehat{\mu}^1_{\text{agg}}(\eta^*) \right) = \sum_{j=1}^K (\widehat{\eta}_j - \eta^*_j) \cdot \sqrt{N_{\mathrm{all}}}(\widehat{\mu}^1_{t,s_j} - \widehat{\mu}^1_{ot})
\]

For $j \in \mathcal{S}^*$: By Lemma \ref{lem:selection}(1), $|\widehat{\eta}_j - \eta^*_j| = O_p(N_{\mathrm{all}}^{-1/2})$. By Theorem \ref{thm:main} with the same rescaling argument as in Lemma \ref{lem:if_agg}, $\sqrt{N_{\mathrm{all}}}(\widehat{\mu}^1_{t,s_j} - \widehat{\mu}^1_{ot}) = O_p(1)$. Thus each term is $o_p(1)$.

For $j \notin \mathcal{S}^*$: By Lemma \ref{lem:selection}(2), $\widehat{\eta}_j = 0$ w.h.p., and $\eta^*_j = 0$ by definition (oracle excludes biased sites). Thus the difference is zero w.h.p.
Because $K$ is fixed, summing over $j$ yields $o_p(1)$.
\end{proof}

\begin{remark}[Growing number of sources $K$]\label{rem:growing_k}
The arguments above use Assumption \ref{asmp:aggregation}(d) (fixed $K$) in two places: (i) to ensure the plug-in term $\sum_{j=1}^K(\widehat{\eta}_j-\eta_j^*)\sqrt{N_{\mathrm{all}}}(\widehat{\mu}^1_{t,s_j}-\widehat{\mu}^1_{ot})$ is $o_p(1)$, and (ii) to treat the variance/covariance estimation problem as finite-dimensional. If $K=K_N\to\infty$, then additional high-dimensional control is needed, e.g., a sufficient condition for Lemma \ref{lem:asymp_equiv} is
\[
\|\widehat{\eta}-\eta^*\|_1 \cdot \max_{1\le j\le K}\left|\sqrt{N_{\mathrm{all}}}\,(\widehat{\mu}^1_{t,s_j}-\widehat{\mu}^1_{ot})\right| = o_p(1),
\]
which typically requires both a sparsity assumption on $\eta^*$ and concentration bounds scaling with $\log K$.
\end{remark}

\subsection{Main Results for Aggregation}

\begin{theorem}[Asymptotic Normality of FACE-HD]\label{thm:face_hd}
Suppose Assumptions \ref{asmp:data}-\ref{asmp:rsc} and \ref{asmp:aggregation} hold. For each informative source $s_j \in \mathcal{S}^*$, let $N_{t,s_j}=N_t+N_{s_j}$ and $s_{\mathrm{tot},j}=s_{\alpha,j}+s_{\gamma,j}$ denote the pairwise total sparsity. If
\[
\max_{j\in\mathcal{S}^*}\frac{s_{\mathrm{tot},j}^2\log^2(p)}{N_{t,s_j}} \to 0,
\]
(equivalently under fixed site ratios, the same condition with $N_{\mathrm{all}}$), then:
\[
\sqrt{N_{\mathrm{all}}}\left( \widehat{\mu}^1_{\text{agg}} - \mu^1 \right) \xrightarrow{d} \mathcal{N}(0, L^*(\eta^*))
\]
where $L^*(\eta^*)$ is the minimum asymptotic variance achieved by the oracle aggregation.

Moreover, the variance estimator is consistent:
\[
\widehat{V} = \widehat{L}(\widehat{\eta}) \xrightarrow{p} L^*(\eta^*)
\]
and the standardized statistic is asymptotically standard normal:
\[
\sqrt{N_{\mathrm{all}} / \widehat{V}} \left( \widehat{\mu}^1_{\text{agg}} - \mu^1 \right) \xrightarrow{d} \mathcal{N}(0, 1)
\]
\end{theorem}

\begin{proof}
Combining Lemmas \ref{lem:if_agg}, \ref{lem:selection}, and \ref{lem:asymp_equiv}:
\[
\sqrt{N_{\mathrm{all}}}(\widehat{\mu}^1_{\text{agg}} - \mu^1) = \sqrt{N_{\mathrm{all}}}(\widehat{\mu}^1_{\text{agg}}(\eta^*) - \mu^1) + o_p(1) = \frac{1}{\sqrt{N_{\mathrm{all}}}} \sum_{i=1}^{N_{\mathrm{all}}} \varphi_{\mathrm{agg},i}(\eta^*) + o_p(1)
\]

As in Lemma \ref{lem:clt}, the observations are independent with a multi-site structure and site proportions converging to constants. The aggregated influence function contributions $\{\varphi_{\mathrm{agg},i}(\eta^*)\}_{i=1}^{N_{\mathrm{all}}}$ are therefore independent with mean zero and finite variance $L^*(\eta^*)$. Under Assumption \ref{asmp:overlap}(b), these contributions have a uniform $(2+\delta)$ moment, so the same Lyapunov argument as in Lemma \ref{lem:clt} applies and yields:
\[
\frac{1}{\sqrt{N_{\mathrm{all}}}} \sum_{i=1}^{N_{\mathrm{all}}} \varphi_{\mathrm{agg},i}(\eta^*) \xrightarrow{d} \mathcal{N}(0, L^*(\eta^*))
\]

For variance estimation consistency: By Lemma \ref{lem:selection}(1), $\|\widehat{\eta} - \eta^*\|_2 = O_p(N_{\mathrm{all}}^{-1/2})$. By the uniform convergence of $\widehat{L}$ to $L^*$ and local Lipschitz continuity of the quadratic map $L^*(\eta)=\eta^T H\eta+g^T\eta+c$ in a neighborhood of $\eta^*$ (e.g., on $\{\eta:\|\eta-\eta^*\|_2\le R\}$,
$|L^*(\eta_1)-L^*(\eta_2)|\le (2\|H\|_{\mathrm{op}}R+\|g\|_2)\|\eta_1-\eta_2\|_2$):
\[
|\widehat{L}(\widehat{\eta}) - L^*(\eta^*)| \le |\widehat{L}(\widehat{\eta}) - L^*(\widehat{\eta})| + |L^*(\widehat{\eta}) - L^*(\eta^*)| = O_p(N_{\mathrm{all}}^{-1/2}) + O_p(N_{\mathrm{all}}^{-1/2}) = o_p(1)
\]

Finally, write
\[
\sqrt{\frac{N_{\mathrm{all}}}{\widehat{V}}}\left(\widehat{\mu}^1_{\text{agg}}-\mu^1\right)
=
\left(\sqrt{\frac{L^*(\eta^*)}{\widehat{V}}}\right)\cdot
\left(\frac{1}{\sqrt{L^*(\eta^*)}}\sqrt{N_{\mathrm{all}}}\left(\widehat{\mu}^1_{\text{agg}}-\mu^1\right)\right).
\]
We already proved $\sqrt{N_{\mathrm{all}}}(\widehat{\mu}^1_{\text{agg}}-\mu^1)\xrightarrow{d}\mathcal{N}(0,L^*(\eta^*))$, hence the second factor converges in distribution to $\mathcal{N}(0,1)$.
By variance consistency, $\widehat{V}\xrightarrow{p}L^*(\eta^*)$, so the prefactor $\sqrt{L^*(\eta^*)/\widehat{V}}\xrightarrow{p}1$.
Therefore the product converges in distribution to $\mathcal{N}(0,1)$ by Slutsky's theorem.
\end{proof}

\begin{corollary}[Asymptotic Normality of the Aggregated ATE]\label{cor:ate_agg}
Define the aggregated ATE estimator as
\[
\widehat{\tau}_{\text{agg}}
\;\coloneqq\;
\widehat{\mu}^1_{\text{agg}}-\widehat{\mu}^0_{\text{agg}},
\]
where $\widehat{\mu}^0_{\text{agg}}$ is constructed by applying the same aggregation procedure to the control-arm estimators $\{\widehat{\mu}^0_{t,s_j}\}_{j=1}^K$ (with the corresponding nuisance estimates and variance/covariance components).
Assume the analogues of Assumptions \ref{asmp:data}--\ref{asmp:aggregation} hold for both arms, so that both $\widehat{\mu}^1_{\text{agg}}$ and $\widehat{\mu}^0_{\text{agg}}$ admit asymptotic linear representations under $\sqrt{N_{\mathrm{all}}}$ scaling.
Then
\[
\sqrt{N_{\mathrm{all}}}\big(\widehat{\tau}_{\text{agg}}-\tau\big)
\;\xrightarrow{d}\;
\mathcal{N}\!\big(0,V_\tau^*\big),
\qquad
V_\tau^* \coloneqq
\Var\!\Big(\varphi_{\mathrm{agg}}^{1}(D)-\varphi_{\mathrm{agg}}^{0}(D)\Big),
\]
where $\varphi_{\mathrm{agg}}^{1}$ and $\varphi_{\mathrm{agg}}^{0}$ denote the aggregated influence functions (oracle-level) for the treated and control aggregated estimators, respectively.
\end{corollary}

\begin{proposition}[Efficiency Gain]\label{prop:efficiency}
Under Assumptions \ref{asmp:data}-\ref{asmp:rsc} and \ref{asmp:aggregation}, the FACE-HD estimator is asymptotically at least as efficient as the target-only estimator:
\[
L^*(\eta^*) \le L^*(0)
\]
with strict inequality whenever there exists $j\in\mathcal{S}^*$ such that
\[
\Var(\varphi_{t,s_j} - \varphi_{ot}) > 0
\quad\text{and}\quad
\Cov(\varphi_{ot}, \varphi_{t,s_j} - \varphi_{ot}) \ne 0.
\]
\end{proposition}

\begin{proof}
Since $\eta^*$ minimizes $L^*(\eta)$ and $\eta = 0$ is feasible, we have $L^*(\eta^*) \le L^*(0)$.

For strict improvement, consider one informative source $j \in \mathcal{S}^*$. The optimal projection is:
\[
\eta^*_j = -\frac{\text{Cov}(\varphi_{ot}, \varphi_{t,s_j} - \varphi_{ot})}{\Var(\varphi_{t,s_j} - \varphi_{ot})}
\]
If the covariance is non-zero and the denominator is positive, the minimizing weight is non-zero and strictly decreases variance relative to $\eta=0$, so $L^*(\eta^*) < L^*(0)$.
\end{proof}

\bibliographystyle{plainnat}
\bibliography{ref}

\end{document}
